\documentclass[10pt,a4paper]{article}
\usepackage[margin=2.2cm]{geometry}
\usepackage{fontspec}
\setmonofont{DejaVu Sans Mono}[Scale=0.84]
\newfontfamily\unifont{DejaVu Serif}[Scale=0.9]
\usepackage{amsmath,amssymb,mathtools}
\setcounter{MaxMatrixCols}{20}
\usepackage{longtable,booktabs,array}
\usepackage{fancyvrb}
\usepackage[dvipsnames]{xcolor}
\usepackage[colorlinks=true,linkcolor=NavyBlue,urlcolor=NavyBlue]{hyperref}
\usepackage{microtype}
\usepackage{needspace}
\usepackage{tocloft}\setlength{\cftbeforesecskip}{0.45em}
\setlength{\parindent}{0pt}\setlength{\parskip}{0.55em}
\setlength{\emergencystretch}{8em}
\makeatletter
\renewcommand\section{\@startsection{section}{1}{\z@}{-3.5ex \@plus -1ex \@minus -.2ex}{2.3ex \@plus.2ex}{\normalfont\Large\bfseries\raggedright}}
\renewcommand\subsection{\@startsection{subsection}{2}{\z@}{-3.25ex\@plus -1ex \@minus -.2ex}{1.5ex \@plus .2ex}{\normalfont\large\bfseries\raggedright}}
\renewcommand\subsubsection{\@startsection{subsubsection}{3}{\z@}{-3.25ex\@plus -1ex \@minus -.2ex}{1.5ex \@plus .2ex}{\normalfont\normalsize\bfseries\raggedright}}
\makeatother
\renewcommand{\arraystretch}{1.15}
\allowdisplaybreaks
\begin{document}
\title{The world brane of 7+1 spacetime\\[0.4em]\large Where the Standard Model lives, how it sits in the bulk, and how it meets gravity and dirac16complex — a guide for students, with worked examples checked in sections 25 and 26 of the notebook 7+1.nb and by independent Python programs}
\author{Repository: github.com/once-ere/7plus1 — prepared for Nash\\ Notebook: 7+1.nb (Wolfram Language 15.0.1), evaluated Wed 28 Oct 2026 08:57:47\\ Document generated by tools/make\_docs.py on 2026-10-28}
\date{}
\maketitle
\begin{abstract}
The WORLD BRANE of the 7+1 model is the 3+1-dimensional submanifold $W=\{x_4=x_5=x_6=x_7=0\}$ of the eight-dimensional spacetime: the place, a single point of the four-dimensional auxiliary torus, where the Standard Model lives, while gravity and the 16-component field dirac16complex fill the whole bulk. This document defines the brane precisely — its embedding, its induced metric $h_{\mu\nu}=g_{\mu\nu}(x,\vec y=0)$ and signature, its symmetries, its extrinsic curvature — and explains, with worked examples, how integrals over it are written in eight dimensions, what energy–momentum it carries (a tension with $w=-1$, brane fluids), how gravity looks from it (the four-dimensional Friedmann equations with $8\pi G_4=\kappa/V_\perp$, the radion coupling, Newton’s law), how the 16 components of dirac16complex appear on it as four brane Dirac spinors (with the explicit matrices), and how the portals couple the Standard Model to every Kaluza–Klein level at once, because the brane breaks the translations of the torus. It follows the brane through the three regimes of the model, and compares it with a brane at the north pole of a 4-SPHERE (the same coupling to the radion, one state per Kaluza–Klein level, no zero modes of dirac16complex) and with the WARPED branes of Randall and Sundrum (tensions, junction conditions, a TeV-strong radion), and shows that the mixing of the Higgs boson with the radion, which the brane allows, is negligible in 7+1. Every equation used is written here. Section 25 of 7+1.nb contains 27 checks (27 True); the independent program tools/ref\_torus\_brane.py contains 38 (38 True); the 17 groups of numbers printed by both agree to a largest relative difference of $3.8\times10^{-10}$. Section 26 (the radion, the sphere and the warped comparison) contains 33 checks (33 True), tools/ref\_radion.py 31 (31 True), and their 15 groups of numbers agree to $8.2\times10^{-11}$.
\end{abstract}
\tableofcontents
\Needspace{12\baselineskip}\section*{1. Results at a glance}\addcontentsline{toc}{section}{1. Results at a glance}
{\footnotesize
\begin{longtable}{>{\raggedright\arraybackslash}p{\dimexpr 0.2400\linewidth-2\tabcolsep\relax}>{\raggedright\arraybackslash}p{\dimexpr 0.6400\linewidth-2\tabcolsep\relax}>{\raggedright\arraybackslash}p{\dimexpr 0.1200\linewidth-2\tabcolsep\relax}}
\toprule
\textbf{question} & \textbf{answer} & \textbf{where} \\
\midrule\endfirsthead
\toprule
\textbf{question} & \textbf{answer} & \textbf{where} \\
\midrule\endhead
\endfoot
\bottomrule\endlastfoot
What is the world brane? & the 3+1-dimensional submanifold $W=\{x_4=x_5=x_6=x_7=0\}$ of the bulk, through the point $\vec y=0$ of the auxiliary torus; spanned by time $x_0$ and ordinary space $x_1,x_2,x_3$ & §3.2 \\
What metric do brane observers use? & the induced metric $h_{\mu\nu}(x)=g_{\mu\nu}(x,0)$; flat: $\mathrm{diag}(+1,-1,-1,-1)$, the squares of the quaternion basis $e_0,e_1,e_2,e_3$; signature list $\{3,1,0\}$ & §3.3 \\
Does the brane see the size of the torus? & not in its metric: for the cosmological metric $h=\mathrm{diag}(1,-a^2,-a^2,-a^2)$; it feels $b$ through $G_4=G_8/(b^4L_c^4)$ and through the radion coupling & §3.3, §7 \\
How does the brane sit in the bulk? & totally geodesic for the metrics of the model (extrinsic curvature $K^{(k)}_{\mu\nu}=\frac1{2b}\partial_kg_{\mu\nu}=0$); example of a metric where it is not: $\{4\,00\}=1458/58619$ & §4 \\
How is a brane integral written in eight dimensions? & $\int_Wd^4x\sqrt{|h|}\,\mathcal L=\int d^8X\sqrt{|g|}\,\frac{\delta^4(\vec y)}{b^4}\,\mathcal L$ & §5 \\
What does a tension do? & $T^W_{\mu\nu}=f_Wh_{\mu\nu}$: energy $f_W$, pressure $-f_W$ ($w=-1$); smeared $\varepsilon=f_W/V_\perp$, $p_a=-\varepsilon$, $p_b=0$; a part of the four-dimensional cosmological constant & §6, §7.1 \\
How does gravity look from the brane? & four-dimensional Friedmann equations with $8\pi G_4=\kappa/V_\perp$; the age of the universe 13.79 Gyr; Newton’s law at $r\gg\ell_\perp$; the radion couples by $\phi\,T^\mu{}_\mu/(\sqrt3M_{\rm Pl})$ & §7 \\
What lives on it? & the Standard Model: $SU(3)\times SU(2)\times U(1)$, three generations with right-handed neutrinos, the Higgs doublet; its gauge fields are confined to the brane (Coulomb’s law at all distances) & §8, §7.4 \\
What does dirac16complex look like there? & four brane Dirac spinors $\psi_a$ ($\psi=U\zeta$, $\zeta=\sum_a\psi_a\otimes e_a$); the brane Lorentz group acts on the spinor index, the rotations about the brane on $a$, which carries $(2,1)\oplus(1,2)$ of $Spin(4)$ & §9 \\
How does the brane couple to the bulk? & by local portals at $\vec y=0$; every Kaluza–Klein level has the same value $1/\sqrt{V_\perp}$ there, so the portals couple the whole tower and do not conserve Kaluza–Klein momentum & §10 \\
What does the brane do in the three regimes? & negligible in the primordial regime ($2.0\times10^{-28}$ of the energy at $T=H=M_8$); hot below the normalcy temperature $T_*=13.3$ GeV; the Friedmann universe today; diluted in the de Sitter future & §11 \\
Does the Higgs boson on the brane mix with the radion? & through $\xi R(h)\Phi^\dagger\Phi$, with the parameter $\gamma=v/(\sqrt3M_{\rm Pl})=5.84\times10^{-17}$: mixing angle $-3.53\times10^{-16}$ for $\xi=1$, $m_{r0}=10$ GeV — negligible & §7.5 \\
What if the extra space were a 4-sphere? & a brane at the north pole has the same induced metric $e^{-4\beta}g_E$ and radion coupling; it keeps $SO(4)$ of $SO(5)$, sees one ZONAL state per level with the strength $\sqrt{d_l/V}$, and gets no massless fermions from dirac16complex & §12 \\
And a warped bulk? & Randall–Sundrum: tensions $\pm6k/\kappa$ (UV positive), radion coupling $1/\langle\mu\rangle$, $1.41\times10^{15}$ times stronger for $\langle\mu\rangle=3$ TeV; in eight dimensions with only $\Lambda$ no warped solution exists & §13 \\*
\end{longtable}}
\Needspace{12\baselineskip}\section*{2. Conventions}\addcontentsline{toc}{section}{2. Conventions}
Units $\hbar=c=1$ (energies and masses in GeV; $\hbar c=1.973269804\times10^{-16}$ GeV m). The coordinates of 7+1 spacetime are $X=(x_0,x_1,\dots,x_7)$: $x_0$ is time (the direction of the octonion unit $e_0=1$), $x_1,x_2,x_3$ are ordinary space, and $x_4,\dots,x_7$ the four auxiliary space-like coordinates, collected in $\vec y=(x_4,\dots,x_7)$; on the brane $x=(x_0,x_1,x_2,x_3)$. Indices: $A,B=0,\dots,7$ (all directions), $\mu,\nu=0,\dots,3$ (along the brane), $k,l=4,\dots,7$ (transverse to it). A prime is $d/dx_0$. The frame metric is $\eta=\mathrm{diag}(+1,-1,\dots,-1)$, simply the squares of the octonion basis (explained in the next paragraph).

\textbf{The signature: $\{7,1,0\}$ and $\mathrm{diag}(+1,-1,\dots,-1)$.} The 7+1 spacetime has one time-like dimension ($x_0$, the direction of the octonion unit $e_0=1$) and seven space-like dimensions ($x_1,x_2,x_3$ ordinary space, $x_4,\dots,x_7$ auxiliary space). Two different objects describe this, and neither may be replaced by the other. (i) The tensor package xAct, used in section 12 of 7+1.nb exactly as in the author’s notebook Einstein-Lovelock-7+1-Nash.nb, declares the metric with \texttt{DefMetric[\{7, 1, 0\}, g71[-a, -b], CD, …]}: the list $\{7,1,0\}$ declares seven space-like dimensions, one time-like dimension and no degenerate dimension — the numbers of entries $+1$ and $-1$ of $\eta_{71}=\mathrm{diag}(-1,+1,\dots,+1)$ (\texttt{η71}), the same spacetime written in the sign convention of 4+4.nb (space-like $=+1$, time-like $=-1$). (ii) The matrix $\eta=-\eta_{71}=\mathrm{diag}(+1,-1,\dots,-1)$ (\texttt{η17}) is simply the squares of the octonion basis: $e_0^2=+1$ for the unit $e_0=1$ (time) and $e_i^2=-1$ for the seven imaginary units $e_1,\dots,e_7$ (space); more generally it is the real part of the product of two basis units, and it gives the quadratic form of the octonions:
\begin{gather*}
\eta_{AB}=\mathrm{Re}(e_Ae_B),\qquad e_0^2=+1,\qquad e_i^2=-1\ \ (i=1,\dots,7),\\ \mathrm{Re}(x\,x)=x_0^2-x_1^2-\cdots-x_7^2=\eta_{AB}\,x^Ax^B\qquad (x=x^Ae_A\in\mathbb O).
\end{gather*}
$\eta$ gives the signs of the Clifford relation $\Gamma_A\Gamma_B+\Gamma_B\Gamma_A=2\,\eta_{AB}\,1_{16}$ of the real 16×16 Dirac matrices: time squares to $+1_{16}$ and each of the seven space directions to $-1_{16}$, exactly like the octonion units (in the notation $Cl(p,q)$, with $p$ generators squaring to $+1$ and $q$ to $-1$, the algebra $Cl(1,7)\cong\mathbb R(16)$). It also gives the signs of the metric components: $g_{00}>0$ along the time-like direction, negative squares along the space-like directions. Every component computed by 7+1.nb follows from the component matrix of the metric; the declared list fixes for xAct the sign of the determinant of the metric, \texttt{SignDetOfMetric[g71]} $=-1$, which is the same in both sign conventions, since $\det \eta=\det \eta_{71}=-1$. 7+1.nb checks \texttt{SignatureOfMetric[g71]} $=\{7,1,0\}$, \texttt{SignDetOfMetric[g71]} $=-1$, $\mathrm{Re}(e_Ae_B)=\eta_{AB}$ and $\mathrm{Re}(x\,x)=\eta_{AB}x^Ax^B$.

\textbf{Physical constants} (as in sections 17–24 of 7+1.nb): $M_8=10$ TeV, $\kappa=8\pi G_8=M_8^{-6}$; $M_{\rm Pl}=2.435\times10^{18}$ GeV, $8\pi G_4=M_{\rm Pl}^{-2}$; the radion mass $m_r=10$ GeV; the Higgs mass $m_h=125.20$ GeV.

\Needspace{12\baselineskip}\section*{3. The world brane: definition}\addcontentsline{toc}{section}{3. The world brane: definition}
\Needspace{8\baselineskip}\subsection*{3.1 The bulk and the auxiliary torus}\addcontentsline{toc}{subsection}{3.1 The bulk and the auxiliary torus}
The four auxiliary coordinates are periodic, $x_k\simeq x_k+L_cn_k$ ($k=4,\dots,7$, $n_k\in\mathbb Z$, $L_c$ the coordinate period): they form the AUXILIARY TORUS $T^4=\mathbb R^4/(L_c\mathbb Z)^4$, a compact, flat space without boundary. The spacetime of the model, the BULK, is
\begin{equation*}
M=\mathbb R_{x_0}\times\mathbb R^3_{\vec x}\times T^4_{\vec y},\qquad b^8=\det(-g_{kl}),\qquad V_\perp=b^4L_c^4,\qquad \ell_\perp=V_\perp^{1/4}=bL_c,
\end{equation*}
with $b$ the size of the torus (for the metrics of the model $-g_{kl}=b^2\delta_{kl}$), $V_\perp$ its proper volume and $\ell_\perp$ the proper circumference of each of its four circles; today $V_\perp=M_{\rm Pl}^2/M_8^6=5.929\times10^{12}$ GeV$^{-4}$ and $\ell_\perp=1560.45$ GeV$^{-1}=3.079\times10^{-13}$ m. The translations $\vec y\to\vec y+\vec c$ are symmetries of the torus, so all its points are alike.

\Needspace{8\baselineskip}\subsection*{3.2 The definition of the world brane}\addcontentsline{toc}{subsection}{3.2 The definition of the world brane}
\textbf{Definition.} The WORLD BRANE is the submanifold
\begin{equation*}
W=\{X\in M:\ x_4=x_5=x_6=x_7=0\}=\mathbb R_{x_0}\times\mathbb R^3_{\vec x}\times\{\vec y=0\}\subset M,
\end{equation*}
the set of points of the bulk that sit at the point $\vec y=0$ of the torus. It is 3+1-dimensional (time and ordinary space), and its codimension is 4 (the four transverse directions $x_4,\dots,x_7$). It is described by its EMBEDDING
\begin{equation*}
X^A(x)=(x_0,\,x_1,\,x_2,\,x_3,\,0,\,0,\,0,\,0),\qquad J^A{}_\mu=\frac{\partial X^A}{\partial x^\mu}=\begin{pmatrix}1_4\\0_{4\times4}\end{pmatrix},
\end{equation*}
with the brane coordinates $x^\mu$ equal to the first four bulk coordinates; the columns of $J$ are the four tangent vectors $\partial_\mu$, and $\partial_4,\dots,\partial_7$ are the normal directions. Because every point of the torus is like every other, the choice $\vec y=0$ is a choice of origin; in 7+1.nb the position of the brane is fixed (it is not a dynamical field).

The Standard Model lives on $W$: its fields are functions of $x$ only. Gravity and dirac16complex live in the bulk and are functions of $x$ and $\vec y$. The total action of 7+1.nb (section 20.1) is the sum of bulk integrals over $d^8X\sqrt{|g|}$ and one brane integral over $d^4x\sqrt{|h|}$ ($h$ the induced metric of section 3.3):
\begin{align*}
S&=S_{\rm grav}+S_\psi+S_U+S_W,\qquad S_{\rm grav}=-\frac1{2\kappa}\int d^8X\sqrt{|g|}\,\big(2\Lambda+R+\alpha_2L_2+\alpha_3L_3\big), \\
S_\psi&=\int d^8X\sqrt{|g|}\,\big[\bar\psi\,i\Gamma^A(X)(\partial_A+\Omega_A)\psi-m\,\bar\psi\psi\big]+S_{\rm self},\qquad S_U=-\int d^8X\sqrt{|g|}\,U(b), \\
S_W&=\int_Wd^4x\sqrt{|h|}\,\big(-f_W+\mathcal L_{\rm SM}+\mathcal L_{\rm portal}\big),
\end{align*}
with $\kappa=8\pi G_8$, $\Lambda$ the bulk cosmological constant, $R$ the scalar curvature, $\alpha_2,\alpha_3$ the Lovelock couplings and
\begin{align*}
L_2&=R^2-4R_{AB}R^{AB}+R_{ABCD}R^{ABCD}\qquad(\text{Gauss–Bonnet}), \\
L_3&=R^3-12RR_{AB}R^{AB}+16R^A{}_BR^B{}_CR^C{}_A+24R^{AC}R^{BD}R_{ABCD}+3RR_{ABCD}R^{ABCD} \\
&\quad-24R^{AB}R_A{}^{CDE}R_{BCDE}+4R^{AB}{}_{CD}R^{CD}{}_{EF}R^{EF}{}_{AB}-8R_A{}^C{}_B{}^DR_C{}^E{}_D{}^FR_E{}^A{}_F{}^B\qquad(\text{cubic Lovelock});
\end{align*}
$\psi$ is dirac16complex with $\bar\psi=\psi^\dagger\Gamma_0$, the frame Dirac matrices $\Gamma^A(X)$ and the spin connection $\Omega_A$; $S_{\rm self}$ are its nine self-couplings (not needed here); $U(b)$ is the stabilizing potential of the torus (section 7.3); $f_W$ is the brane tension (section 6), $\mathcal L_{\rm SM}$ the Standard Model (section 8) and $\mathcal L_{\rm portal}$ the couplings between the Standard Model and dirac16complex (section 10). There is no curvature term of the brane itself: brane observers feel gravity only through $h_{\mu\nu}$.

\Needspace{8\baselineskip}\subsection*{3.3 The induced metric}\addcontentsline{toc}{subsection}{3.3 The induced metric}
Brane observers measure lengths and times with the INDUCED METRIC, the bulk metric restricted to the tangent directions of $W$:
\begin{equation*}
h_{\mu\nu}(x)=g_{AB}\big(X(x)\big)\,J^A{}_\mu J^B{}_\nu=g_{\mu\nu}(x,\vec y=0),\qquad h=J^{\mathsf T}gJ .
\end{equation*}
For flat spacetime $h=\mathrm{diag}(+1,-1,-1,-1)$: one time-like and three space-like dimensions; in the counting of xAct this is the signature list $\{3,1,0\}$ (three space-like, one time-like and no degenerate dimension; the bulk is declared with $\{7,1,0\}$). These signs are $\eta$ restricted to $A,B=0,\dots,3$, i.e. the squares of the basis $e_0=1$, $e_1,e_2,e_3$ of the QUATERNION SUBALGEBRA $\mathbb H$ of the octonions. In the octonion table of 7+1.nb
\begin{equation*}
e_1e_2=-e_3,\qquad e_2e_3=-e_1,\qquad e_3e_1=-e_2,\qquad e_i^2=-1,\qquad e_ie_j=-\delta_{ij}-\epsilon_{ijk}e_k\quad(i,j,k=1,2,3),
\end{equation*}
so the products of $e_0,\dots,e_3$ stay in their span (the subalgebra is closed) and $\mathrm{Re}(e_\mu e_\nu)=\mathrm{diag}(1,-1,-1,-1)$ (checked in section 17 of 7+1.nb and by the Python program, which reads the coefficients $-1,-1,-1$ of $e_3,e_1,e_2$ in $e_1e_2,e_2e_3,e_3e_1$ off the octonion table).

\Needspace{6\baselineskip}\subsubsection*{Example 3.1 — the cosmological metric: an expanding universe on the brane}
For the cosmological metric of section 11.2 of 7+1.nb, the metric of all three regimes of the model (sections 20–23),
\begin{equation*}
ds^2=dx_0^2-a^2(x_0)\,(dx_1^2+dx_2^2+dx_3^2)-b^2(x_0)\,(dx_4^2+\cdots+dx_7^2)
\end{equation*}
the induced metric is the Friedmann–Robertson–Walker metric
\begin{equation*}
h=J^{\mathsf T}gJ=\mathrm{diag}(1,-a^2,-a^2,-a^2),\qquad ds_W^2=dx_0^2-a^2(x_0)\,d\vec x^2 :
\end{equation*}
brane observers see an expanding universe with the scale factor $a$, and nothing of $b$ (section 25.6 of 7+1.nb).

\Needspace{6\baselineskip}\subsubsection*{Example 3.2 — an inhomogeneous metric}
Section 11.3 of 7+1.nb tests its curvature engine with the metric $ds^2=n^2dx_0^2-a^2d\vec x^2-b^2d\vec y^2$ with the polynomials
\begin{equation*}
n=1+\frac{x_0x_4}5+\frac{x_0^2}{11},\qquad a=1+\frac{x_0}3+\frac{x_4^2}7,\qquad b=2-\frac{x_0x_4}4+\frac{x_0^2}9+\frac{x_4}{13},
\end{equation*}
which depend on time AND on the transverse coordinate $x_4$. At the brane ($x_4=0$) the induced metric is
\begin{equation*}
h=\mathrm{diag}\Big(\big(1+\tfrac{x_0^2}{11}\big)^2,\ -\big(1+\tfrac{x_0}3\big)^2,\ -\big(1+\tfrac{x_0}3\big)^2,\ -\big(1+\tfrac{x_0}3\big)^2\Big):
\end{equation*}
the brane sees only the values at $x_4=0$; how the metric changes off the brane enters its extrinsic curvature (section 4).

\Needspace{6\baselineskip}\subsubsection*{Example 3.3 — the primordial regime seen from the brane}
In the primordial regime of the model the cubic, the Gauss–Bonnet and finally the Einstein term of gravity dominate one after another (section 21 of 7+1.nb). In its Einstein part the vacuum metric is the Kasner solution $a\propto\tau^{p_1}$, $b\propto\tau^{q_1}$ with
\begin{equation*}
p_1=\frac{1+2\sqrt2}7=0.546918,\qquad q_1=\frac{2-3\sqrt2}{14}=-0.160189,\qquad 3p_1+4q_1=1,\qquad 3p_1^2+4q_1^2=1
\end{equation*}
($\tau$ the time since the singularity; section 21.1). The brane sees the power-law universe $ds_W^2=d\tau^2-\tau^{2p_1}d\vec x^2$, which expands a little faster than a radiation-dominated ($a\propto\tau^{1/2}$) and more slowly than a matter-dominated ($a\propto\tau^{2/3}$) four-dimensional universe, while the torus shrinks ($q_1<0$). The brane does not see the shrinking in its metric; it sees it only through $8\pi G_4=\kappa/V_\perp$, which grows as $V_\perp\propto b^4$ falls.

\Needspace{8\baselineskip}\subsection*{3.4 Symmetries of the brane}\addcontentsline{toc}{subsection}{3.4 Symmetries of the brane}
\begin{itemize}
\item \textbf{Along the brane:} the Lorentz group $SO(1,3)$ of the flat brane (and the translations in $x$), generated in the spinor representation by $\Gamma_\mu\Gamma_\nu$ (section 9).
\item \textbf{About the brane:} the rotations of the four transverse directions, $SO(4)$, generated by $\Gamma_k\Gamma_l$. Locally the bulk near the brane is symmetric under them; on the torus only the $2^4\cdot4!=384$ signed permutations of the axes remain (they fix $\vec y=0$), among them the reflections $y_k\to-y_k$.
\item \textbf{Broken by the brane:} the translations of the torus (the brane is at one point; section 10) and the rotations that would turn the brane into the transverse directions (generated by $\Gamma_\mu\Gamma_k$). The Lorentz group $SO(1,7)$ of the bulk is reduced to $SO(1,3)\times SO(4)$.
\end{itemize}
\Needspace{12\baselineskip}\section*{4. How the brane sits in the bulk: extrinsic curvature}\addcontentsline{toc}{section}{4. How the brane sits in the bulk: extrinsic curvature}
The induced metric says how the brane is curved within itself. How it is bent within the bulk is measured by its EXTRINSIC CURVATURE, one symmetric tensor for each of the four unit normals $n_{(k)}=b^{-1}\partial_k$: half the Lie derivative of the metric along the normal, $K^{(k)}_{\mu\nu}=\frac12(\mathcal L_{n_{(k)}}g)_{\mu\nu}$, i.e. how fast the induced metric changes when the brane is moved along $n_{(k)}$. For metrics with $g_{\mu k}=0$ and $-g_{kl}=b^2\delta_{kl}$ at the brane (all metrics of this document)
\begin{equation*}
K^{(k)}_{\mu\nu}=\frac1{2b}\,\partial_kg_{\mu\nu}\Big|_W,\qquad \{k\,\mu\nu\}\Big|_W=\frac1{2b^2}\,\partial_kg_{\mu\nu}\Big|_W=\frac1b\,K^{(k)}_{\mu\nu},
\end{equation*}
where $\{\lambda\,\mu\nu\}=\frac12g^{\lambda\sigma}(\partial_\mu g_{\sigma\nu}+\partial_\nu g_{\sigma\mu}-\partial_\sigma g_{\mu\nu})$ are the Christoffel symbols. They appear in the geodesic equation of a free particle of the bulk (a dot is $d/ds$, $s$ the proper time):
\begin{equation*}
\ddot y^k=-\{k\,\mu\nu\}\,\dot x^\mu\dot x^\nu-2\{k\,\mu l\}\,\dot x^\mu\dot y^l-\{k\,lm\}\,\dot y^l\dot y^m .
\end{equation*}
For a particle that moves along the brane ($\vec y=0$ and $\dot{\vec y}=0$ at the moment considered) only the first term remains:
\begin{equation*}
\ddot y^k=-\{k\,\mu\nu\}\,\dot x^\mu\dot x^\nu=-\frac1{2b^2}\,\partial_kg_{\mu\nu}\,\dot x^\mu\dot x^\nu .
\end{equation*}
If all $K^{(k)}_{\mu\nu}$ vanish, the brane is TOTALLY GEODESIC: a free particle that moves along it stays on it.

\Needspace{6\baselineskip}\subsubsection*{Example 4.1 — the cosmological metric: the brane is totally geodesic}
The metric of Example 3.1 does not depend on $\vec y$, so $\partial_kg_{\mu\nu}=0$ and every $\{k\,\mu\nu\}$ ($k=4,\dots,7$, $\mu,\nu=0,\dots,3$) vanishes at the brane (section 25.6 checks all 64 symbols). More generally, any metric that is even under the reflections $y_k\to-y_k$ (which fix the brane) has $\partial_kg_{\mu\nu}=0$ at $y_k=0$, because the derivative of an even function vanishes at the origin: a brane at a fixed point of such reflections is automatically totally geodesic.

\Needspace{6\baselineskip}\subsubsection*{Example 4.2 — an inhomogeneous metric: a brane that pushes particles off}
For the metric of Example 3.2 at $x_0=\frac12$ the only non-zero transverse Christoffel symbol at the brane is
\begin{equation*}
\{4\,00\}=\frac{n\,\partial_4n}{b^2}\Big|_{x_0=1/2,\,x_4=0}=\frac{\frac{45}{44}\cdot\frac1{10}}{\big(\frac{73}{36}\big)^2}=\frac{1458}{58619}=0.02487248,
\end{equation*}
($n=\frac{45}{44}$, $\partial_4n=x_0/5=\frac1{10}$, $b=\frac{73}{36}$ there; $\{4\,ii\}=-a\,\partial_4a/b^2=0$ because $\partial_4a=2x_4/7$ vanishes at the brane). A free particle momentarily at rest on the brane ($dx_i/dx_0=0$, $dx_k/dx_0=0$) therefore has, in coordinate time, $d^2x_4/dx_0^2=-\{4\,00\}=-1458/58619$: it is pushed off the brane, towards negative $x_4$, because this metric is not even in $x_4$ (section 25.6 of 7+1.nb; exact rational arithmetic in both programs).

\Needspace{8\baselineskip}\subsection*{4.1 A thin brane}\addcontentsline{toc}{subsection}{4.1 A thin brane}
The brane of 7+1.nb is infinitely thin (a delta function in $\vec y$). That is an idealization: the linearized gravitational field of its tension falls off like the Green’s function of the four transverse directions,
\begin{equation*}
\nabla_{\vec y}^2\frac1{|\vec y|^2}=-4\pi^2\,\delta^4(\vec y),\qquad G_\perp(\vec y)=-\frac1{4\pi^2|\vec y|^2},\qquad |\delta g|\sim\frac{\kappa f_W}{4\pi^2|\vec y|^2},
\end{equation*}
which is of order one within $|\vec y|\approx(\kappa f_W)^{1/2}/(2\pi)$ of the brane: for $f_W=M_8^4$ that is $1/(2\pi M_8)=1.59\times10^{-5}$ GeV$^{-1}$ (ESTIMATE). There the thin-brane description fails; at distances large compared with $1/M_8$ (and $\ell_\perp=1560.45$ GeV$^{-1}$ is about $10^8$ times larger) the smeared description of sections 5–6 holds.

\Needspace{12\baselineskip}\section*{5. Integrals on the brane, and the delta function of the torus}\addcontentsline{toc}{section}{5. Integrals on the brane, and the delta function of the torus}
A brane action is an integral over $W$ with the volume element $\sqrt{|h|}\,d^4x$. Since the metrics of the model are block diagonal at the brane, $\sqrt{|g|}=\sqrt{|h|}\,\sqrt{\det(-g_{kl})}=\sqrt{|h|}\,b^4$ there, and the brane integral can be written as an integral over the whole bulk:
\begin{equation*}
\int_Wd^4x\,\sqrt{|h|}\;\mathcal L(x)=\int d^8X\,\sqrt{|g|}\;\frac{\delta^4(\vec y)}{b^4}\;\mathcal L(x),\qquad \delta^4(\vec y)=\delta(x_4)\delta(x_5)\delta(x_6)\delta(x_7).
\end{equation*}
\Needspace{6\baselineskip}\subsubsection*{Example 5.1 — the two volume elements, and a test function}
For the cosmological metric $\sqrt{|g|}=a^3b^4$ and $\sqrt{|h|}=a^3$; for the inhomogeneous metric at $x_0=\frac12$, $x_4=0$: $\sqrt{|g|}=na^3b^4$, $\sqrt{|h|}=na^3$, $b^4=(73/36)^4$, exactly. And for any test function $F$,
\begin{equation*}
\int d^4y\;b^4\,\frac{\delta^4(\vec y)}{b^4}\,F(\vec y)=F(0),\qquad\text{e.g. }F=e^{x_4+2x_5}\cos(x_6-x_7):\ F(0)=1
\end{equation*}
(section 25.6 of 7+1.nb). On the torus the delta function is periodic and contains every Fourier mode with the same weight,
\begin{equation*}
\delta^4(\vec y)=\frac1{L_c^4}\sum_{\vec n\in\mathbb Z^4}e^{i\vec q_{\vec n}\cdot\vec y},\qquad \vec q_{\vec n}=\frac{2\pi\vec n}{L_c},\qquad \int_{-L_c/2}^{L_c/2}e^{-iq_ny}\,\delta(y)\,dy=1 ;
\end{equation*}
its average over the torus (its zero mode) is $1/L_c^4$. A brane source of four-dimensional energy density $\rho_4$ has the eight-dimensional density $\rho_4\delta^4(\vec y)/b^4$, whose average over the torus is $\rho_4/V_\perp$ — the SMEARED source of the homogeneous equations (section 25.5 of 7+1.nb).

\Needspace{12\baselineskip}\section*{6. Matter on the brane: the energy–momentum tensor}\addcontentsline{toc}{section}{6. Matter on the brane: the energy–momentum tensor}
A brane action $S_W=\int_Wd^4x\sqrt{|h|}\,\mathcal L_W$ gives the brane energy–momentum tensor $T^W_{\mu\nu}=(2/\sqrt{|h|})\,\delta S_W/\delta h^{\mu\nu}$ (the convention of 7+1.nb: $T^0{}_0$ is the energy density), and the eight-dimensional source
\begin{equation*}
T^{AB}(x,\vec y)=\frac{\delta^4(\vec y)}{b^4}\;T_W^{\mu\nu}(x)\,\delta^A{}_\mu\,\delta^B{}_\nu :
\end{equation*}
it has NO components along the torus — brane matter does not push on the transverse directions.

\Needspace{6\baselineskip}\subsubsection*{Example 6.1 — the tension}
For $\mathcal L_W=-f_W$ (a constant energy per brane volume) use $\delta\sqrt{|h|}=-\frac12\sqrt{|h|}\,h_{\mu\nu}\delta h^{\mu\nu}$:
\begin{equation*}
T^W_{\mu\nu}=\frac2{\sqrt{|h|}}\,\frac{\delta}{\delta h^{\mu\nu}}\Big(-f_W\sqrt{|h|}\Big)=f_W\,h_{\mu\nu},\qquad T^0{}_0=T^1{}_1=T^2{}_2=T^3{}_3=f_W :
\end{equation*}
the energy density is $\varepsilon=T^0{}_0=f_W$ and the pressure along the brane $p=-T^1{}_1=-f_W$ (for a perfect fluid $T^\mu{}_\nu=\mathrm{diag}(\varepsilon,-p,-p,-p)$), so $w=p/\varepsilon=-1$, like a cosmological constant on the brane. Section 25.8 checks this for the diagonal metric $h=\mathrm{diag}(1/u_0,-1/u_1,-1/u_2,-1/u_3)$ by differentiating $-f_W/\sqrt{u_0u_1u_2u_3}$.

\Needspace{8\baselineskip}\subsection*{6.1 Smeared sources in the homogeneous equations}\addcontentsline{toc}{subsection}{6.1 Smeared sources in the homogeneous equations}
With the homogeneous metric
\begin{equation*}
ds^2=n^2dx_0^2-a^2\,d\vec x^2-b^2\,d\vec y^2\qquad(n,\ a,\ b\ \text{functions of }x_0)
\end{equation*}
every source has the form
\begin{equation*}
T^A{}_B=\mathrm{diag}(\varepsilon,-p_a,-p_a,-p_a,-p_b,-p_b,-p_b,-p_b),
\end{equation*}
with the energy density $\varepsilon$, the pressure $p_a$ along ordinary space and the pressure $p_b$ along the torus. Let $W(a,b)$ be the energy in a cell of unit coordinate volume ($\Delta x_1\cdots\Delta x_7=1$, proper volume $V=a^3b^4$) at fixed occupation numbers. The first law $dW=-p_a\,(3a^2b^4\,da)-p_b\,(4a^3b^3\,db)$ gives
\begin{align*}
\varepsilon&=\frac W{a^3b^4},\qquad p_a=-\frac{\partial_aW}{3a^2b^4},\qquad p_b=-\frac{\partial_bW}{4a^3b^3}, \\
0&=\varepsilon'+3H_a(\varepsilon+p_a)+4H_b(\varepsilon+p_b)\qquad(H_a=a'/a,\ H_b=b'/b;\ \text{conservation of energy, for every }W).
\end{align*}
A brane source, smeared over the torus ($\delta^4(\vec y)\to L_c^{-4}$), has a comoving energy that does not depend on $b$, because a brane cell has the volume $a^3$ only:
\begin{align*}
\text{tension: }W&=\frac{a^3f_W}{L_c^4}\ \Rightarrow\ \varepsilon=\frac{f_W}{V_\perp},\quad p_a=-\frac{f_W}{V_\perp},\quad p_b=0, \\
\text{brane fluid, }\rho_4\propto a^{-3(1+w)}:\ W&=\frac{a^3\rho_4}{L_c^4}\ \Rightarrow\ \varepsilon=\frac{\rho_4}{V_\perp},\quad p_a=w\,\varepsilon,\quad p_b=0 .
\end{align*}
\Needspace{6\baselineskip}\subsubsection*{Example 6.2 — radiation on the brane and radiation in the bulk}
Radiation on the brane ($w=\frac13$) has $p_a=\varepsilon/3$ and $p_b=0$. Radiation in the bulk, isotropic in the seven space directions, is different: a massless particle of comoving momenta $k_a$ (ordinary space) and $k_b$ (torus) has $W=\omega=\sqrt{k_a^2/a^2+k_b^2/b^2}$, hence $p_aV=k_a^2/(3a^2\omega)$, $p_bV=k_b^2/(4b^2\omega)$ ($V=a^3b^4$). Write the physical momenta as $k_a/a=p\cos\theta$ and $k_b/b=p\sin\theta$: then $\omega=p$, $p_a/\varepsilon=\cos^2\theta/3$ and $p_b/\varepsilon=\sin^2\theta/4$. For momenta distributed isotropically in seven dimensions the angle $\theta$ ($0\le\theta\le\pi/2$) has the weight $\cos^2\theta\sin^3\theta$ (from the volume element $p_{(3)}^2dp_{(3)}\,p_{(4)}^3dp_{(4)}$ of the three- and four-dimensional parts of the momentum), so
\begin{equation*}
\langle\cos^2\theta\rangle=\frac{\int_0^{\pi/2}\cos^4\theta\sin^3\theta\,d\theta}{\int_0^{\pi/2}\cos^2\theta\sin^3\theta\,d\theta}=\frac{2/35}{2/15}=\frac37,\qquad p_a=\frac13\cdot\frac37\,\varepsilon=\frac\varepsilon7,\qquad p_b=\frac14\cdot\frac47\,\varepsilon=\frac\varepsilon7 .
\end{equation*}
(section 25.8 of 7+1.nb). The equation of state tells whether matter lives on the brane ($p_b=0$) or in the bulk ($p_b\ne0$).

\Needspace{12\baselineskip}\section*{7. Gravity seen from the brane}\addcontentsline{toc}{section}{7. Gravity seen from the brane}
\Needspace{8\baselineskip}\subsection*{7.1 The Friedmann equations of the brane}\addcontentsline{toc}{subsection}{7.1 The Friedmann equations of the brane}
With the torus static ($b=b_0$, $H_b=0$) the field equations of the model reduce, for brane matter, to the four-dimensional Friedmann equations
\begin{equation*}
3H_a^2=8\pi G_4\,\rho_4+\Lambda_4,\qquad H_a'=-\frac\kappa2(\varepsilon+p_a),\qquad 8\pi G_4=\frac{\kappa}{V_\perp}=\frac1{M_{\rm Pl}^2},\qquad \Lambda_4=\Lambda+\kappa U_0+\frac{\kappa f_W}{V_\perp},
\end{equation*}
($H_a=a'/a$; $\Lambda$ the bulk cosmological constant; $U_0$ the value of the stabilizing potential at its minimum; checked in section 20.3 of 7+1.nb). For a brane fluid of energy density $\rho_4$ and pressure $p_4$ ($\varepsilon=\rho_4/V_\perp$, $p_a=p_4/V_\perp$, section 6.1) the second equation is the familiar $H_a'=-4\pi G_4(\rho_4+p_4)$; the tension and the stabilizing potential have $\varepsilon+p_a=0$ and drop out of it. The tension is a part of the four-dimensional cosmological constant.

\Needspace{6\baselineskip}\subsubsection*{Example 7.1 — the age of the universe}
With the Planck 2018 parameters ($H_0=67.4$ km/s/Mpc, $\Omega_m=0.315$, $\Omega_r=9.14\times10^{-5}$, $\Omega_\Lambda=1-\Omega_m-\Omega_r$),
\begin{equation*}
t_0=\frac1{H_0}\int_0^1\frac{da}{a\sqrt{\Omega_ra^{-4}+\Omega_ma^{-3}+\Omega_\Lambda}}=13.791\ \text{Gyr}\qquad(\text{Planck 2018: }13.787\pm0.020)
\end{equation*}
(section 17.5 of 7+1.nb).

\Needspace{6\baselineskip}\subsubsection*{Example 7.2 — the tension as vacuum energy}
For the natural size $f_W=M_8^4=10^{16}$ GeV$^4$ the tension adds to $\Lambda_4$ the vacuum energy $f_W$, which is $3.97\times10^{62}$ times the observed dark-energy density $\rho_\Lambda=3H_0^2M_{\rm Pl}^2\Omega_\Lambda=2.518\times10^{-47}$ GeV$^4$: the other terms of $\Lambda_4$ must cancel it to 62 decimal places — the cosmological-constant problem, which the model does not solve.

\Needspace{8\baselineskip}\subsection*{7.2 Newton’s law between masses on the brane}\addcontentsline{toc}{subsection}{7.2 Newton’s law between masses on the brane}
In linearized eight-dimensional gravity (without the stabilizing potential) a mass $M$ on the brane has the potential $\Phi_8(r)=-5\kappa M/(32\pi^3r^5)$ at short distances. On the torus its images at the lattice points $\ell_\perp\vec n$ add up, and by Poisson summation the sum over images is a sum over the Kaluza–Klein modes:
\begin{align*}
\Phi_{\rm N}(r)&=-\frac{5\kappa M}{32\pi^3}\sum_{\vec n\in\mathbb Z^4}\frac1{(r^2+\ell_\perp^2|\vec n|^2)^{5/2}}=-\frac53\,\frac{G_4M}r\sum_{\vec m\in\mathbb Z^4}e^{-2\pi|\vec m|r/\ell_\perp}, \\
\sum_{\vec n\in\mathbb Z^4}\frac1{(r^2+|\vec n|^2)^{5/2}}&=\frac{4\pi^2}{3r}\sum_{\vec m\in\mathbb Z^4}e^{-2\pi r|\vec m|}\qquad(\text{lengths in units of }\ell_\perp).
\end{align*}
At $r\gg\ell_\perp$ only the zero mode survives, $\Phi_{\rm N}\to-\frac53G_4M/r$: the four-dimensional graviton ($1$) plus the massless radion ($\frac23$, from its coupling below). With the stabilizing potential the radion has the mass $m_r$, its potential falls like $e^{-m_rr}$, and at $r\gg\ell_\perp,1/m_r$ brane observers find Newton’s law $V=-G_4MM'/r$ with $G_4=G_8/V_\perp$. Section 25.4 of 7+1.nb checks the identity to $10^{-8}$ at $r/\ell_\perp=\frac14,\frac12,1,2$.

\Needspace{8\baselineskip}\subsection*{7.3 The radion and brane fields}\addcontentsline{toc}{subsection}{7.3 The radion and brane fields}
The size of the torus is a field, the radion $\beta=\ln(b/b_0)$, held by the stabilizing potential $U(b)=U_0+\frac{K_U}2\ln^2(b/b_0)$ ($m_r^2=\kappa K_U/12$). Brane fields see the metric $h=(b_0/b)^4g_E$, with $g_E$ the four-dimensional Einstein-frame metric; for the canonical radion $\phi=\sqrt{12}M_{\rm Pl}\beta$ this is the coupling
\begin{equation*}
\mathcal L_{\phi W}=\frac\phi{\sqrt3\,M_{\rm Pl}}\,T^\mu{}_\mu,\qquad V_\phi(r)=-\frac23\,\frac{G_4MM'}{r}\,e^{-m_rr}\ \ \text{between two static masses},
\end{equation*}
($T^\mu{}_\mu$ the trace of the energy–momentum tensor of the brane fields): massless radiation ($T^\mu{}_\mu=0$) does not couple, massive particles do. \textbf{Example 7.3:} a radion of 10 GeV decays into brane fields ($b$, $c$, $\tau$, $\mu$ pairs and two gluons) at the rate $\Gamma_{\rm brane}=4.6\times10^{-37}$ GeV (lifetime $1.4\times10^{12}$ s), but into pairs of Kaluza–Klein modes of dirac16complex in the bulk at $\Gamma_{\rm bulk}=6.0\times10^{-24}$ GeV, $1.3\times10^{13}$ times faster (section 22.2 of 7+1.nb).

\Needspace{8\baselineskip}\subsection*{7.4 Gauge forces stay on the brane}\addcontentsline{toc}{subsection}{7.4 Gauge forces stay on the brane}
The photon and the gluons are fields of the brane: their field lines cannot leave it. The static potential of a charge $Q$ on the brane solves the three-dimensional Poisson equation, $V(r)=Q\,Q'/(4\pi r)$ at ALL distances (in Heaviside–Lorentz units), while gravity, a bulk field, changes from $1/r^5$ to $1/r$ at $r\approx\ell_\perp$ (section 7.2). This is why large extra dimensions are compatible with the precise tests of Coulomb’s law and are searched for in gravity and at colliders.

\Needspace{8\baselineskip}\subsection*{7.5 The Higgs boson and the radion: mixing}\addcontentsline{toc}{subsection}{7.5 The Higgs boson and the radion: mixing}
The Higgs doublet $\Phi$ of section 8 is a brane field and may couple to the curvature $R(h)$ of the induced metric. In unitary gauge $\Phi^\dagger\Phi=\varphi_H^2/2$ with $\varphi_H=v+h$ ($v=246$ GeV; here $h$ is the Higgs FIELD, not the metric $h_{\mu\nu}$), and with the canonical radion $\phi$ of section 7.3, given the mass $m_{r0}$ by the stabilizer, the brane Lagrangian of the Higgs boson in the signature $(+,-,-,-)$ is
\begin{equation*}
\mathcal L=\sqrt{|h|}\Big[(D_\mu\Phi)^\dagger D^\mu\Phi-\lambda\Big(\Phi^\dagger\Phi-\frac{v^2}2\Big)^2+\xi\,R(h)\,\Phi^\dagger\Phi\Big]+\frac12(\partial\phi)^2-\frac12m_{r0}^2\phi^2\qquad(\xi=\tfrac16:\ \text{conformal}).
\end{equation*}
On the brane $h_{\mu\nu}=\Omega^2g_{E\mu\nu}$ with $\Omega=e^{-2\beta}=e^{-\phi/(\sqrt3M_{\rm Pl})}$, and for $h_{\mu\nu}=\Omega^2\eta_{\mu\nu}$ the identity $\sqrt{|h|}R(h)=-6\Omega\Box\Omega$ ($\Box=\partial_0^2-\nabla^2$) turns the curvature term into a kinetic mixing. To second order in $(h,\phi)$ and up to total derivatives (section 26.6 of 7+1.nb), with $\gamma=v/(\sqrt3M_{\rm Pl})$ (its sign absorbed into $\phi\to-\phi$):
\begin{equation*}
\mathcal L^{(2)}=\frac12(\partial h)^2+\frac12\big(1+6\xi\gamma^2\big)(\partial\phi)^2+6\xi\gamma\,\partial h\cdot\partial\phi-\frac12m_{h0}^2h^2-\frac12m_{r0}^2\phi^2,\qquad m_{h0}^2=2\lambda v^2 .
\end{equation*}
The shift $h=h'-\frac{6\xi\gamma}Z\phi'$, $\phi=\phi'/Z$ with $Z^2=1+6\xi\gamma^2(1-6\xi)$ makes the kinetic terms canonical, and the rotation $h'=\cos\theta\,h_m-\sin\theta\,\phi_m$, $\phi'=\sin\theta\,h_m+\cos\theta\,\phi_m$ by the angle
\begin{equation*}
\tan2\theta=\frac{12\gamma\xi Z\,m_{h0}^2}{m_{r0}^2-m_{h0}^2\big(Z^2-36\xi^2\gamma^2\big)}
\end{equation*}
diagonalizes the masses ($h_m$, $\phi_m$ the mass eigenstates). \textbf{Example 7.5:} $\gamma=5.838\times10^{-17}$; for $\xi=1$, $m_{r0}=10$ GeV, $m_{h0}=125.20$ GeV the angle is $\theta=-3.525\times10^{-16}$, and the rates of $h_m\to\gamma\gamma$ and $h_m\to gg$ change by $9.95\times10^{-32}$ and $3.39\times10^{-31}$ (60-digit arithmetic in section 26.6 and in tools/ref\_radion.py). The brane allows the mixing, but the flat radion of 7+1, coupled with $1/(\sqrt3M_{\rm Pl})$, cannot be seen through the Higgs boson. (In warped models the coupling scale is $\langle\mu\rangle\sim$ TeV and $\gamma=v/(3\,{\rm TeV})=0.0821$: section 13.)

\Needspace{12\baselineskip}\section*{8. The Standard Model on the brane}\addcontentsline{toc}{section}{8. The Standard Model on the brane}
The Standard Model of 7+1.nb (section 17) has the gauge group $SU(3)_C\times SU(2)_L\times U(1)_Y$ with the gauge fields $G^a_\mu$ ($a=1,\dots,8$), $W^I_\mu$ ($I=1,2,3$), $B_\mu$ and the couplings $g_s,g,g'$, and in each of three generations the Weyl fields (charge $Q=T_3+Y$):

{\footnotesize
\begin{longtable}{>{\raggedright\arraybackslash}p{\dimexpr 0.1400\linewidth-2\tabcolsep\relax}>{\raggedright\arraybackslash}p{\dimexpr 0.1200\linewidth-2\tabcolsep\relax}>{\raggedright\arraybackslash}p{\dimexpr 0.1200\linewidth-2\tabcolsep\relax}>{\raggedright\arraybackslash}p{\dimexpr 0.1200\linewidth-2\tabcolsep\relax}>{\raggedright\arraybackslash}p{\dimexpr 0.5000\linewidth-2\tabcolsep\relax}}
\toprule
\textbf{field} & \textbf{$SU(3)$} & \textbf{$SU(2)$} & \textbf{$Y$} & \textbf{contents} \\
\midrule\endfirsthead
\toprule
\textbf{field} & \textbf{$SU(3)$} & \textbf{$SU(2)$} & \textbf{$Y$} & \textbf{contents} \\
\midrule\endhead
\endfoot
\bottomrule\endlastfoot
$Q_L$ & 3 & 2 & $1/6$ & left-handed quarks $(u_L,d_L)$ \\
$u_R$ & 3 & 1 & $2/3$ & right-handed up-type quarks \\
$d_R$ & 3 & 1 & $-1/3$ & right-handed down-type quarks \\
$L_L$ & 1 & 2 & $-1/2$ & left-handed leptons $(\nu_L,e_L)$ \\
$e_R$ & 1 & 1 & $-1$ & right-handed charged leptons \\
$\nu_R$ & 1 & 1 & $0$ & right-handed neutrinos \\
$\Phi$ & 1 & 2 & $1/2$ & Higgs doublet, $\tilde\Phi=i\sigma_2\Phi^*$ ($\sigma_2$ of section 9.1) \\*
\end{longtable}}
The gauge anomalies cancel in each generation (with or without $\nu_R$), and the relativistic degrees of freedom without $\nu_R$ are $g_*=28+\frac78\cdot90=106.75$ (checked in section 17.1). On the brane the fields feel gravity through $h_{\mu\nu}$, its tetrad and its spin connection; with the $4\times4$ Dirac matrices $\gamma^\mu(x)$ of the brane ($\{\gamma^\mu,\gamma^\nu\}=2h^{\mu\nu}$) and
\begin{equation*}
D_\mu=\nabla_\mu-ig_sG^a_\mu T^a-ig\,W^I_\mu\frac{\tau^I}2-ig'\,Y\,B_\mu
\end{equation*}
($T^a$ the generators of $SU(3)$, $\tau^I=\sigma_I$ the Pauli matrices of section 9.1, $\nabla_\mu$ the covariant derivative of $h$) the Lagrangian is
\begin{align*}
\mathcal L_{\rm SM}&=-\frac14G^a_{\mu\nu}G^{a\mu\nu}-\frac14W^I_{\mu\nu}W^{I\mu\nu}-\frac14B_{\mu\nu}B^{\mu\nu}+(D_\mu\Phi)^\dagger(D^\mu\Phi)-\mu^2\Phi^\dagger\Phi-\lambda(\Phi^\dagger\Phi)^2 \\
&\quad+\sum_f\bar f\,i\gamma^\mu D_\mu f-\big[\bar Q_LY_u\tilde\Phi u_R+\bar Q_LY_d\Phi d_R+\bar L_LY_e\Phi e_R+\bar L_LY_\nu\tilde\Phi\nu_R+\text{h.c.}\big],
\end{align*}
with $G^a_{\mu\nu}=\partial_\mu G^a_\nu-\partial_\nu G^a_\mu+g_sf^{abc}G^b_\mu G^c_\nu$, $W^I_{\mu\nu}=\partial_\mu W^I_\nu-\partial_\nu W^I_\mu+g\epsilon^{IJK}W^J_\mu W^K_\nu$, $B_{\mu\nu}=\partial_\mu B_\nu-\partial_\nu B_\mu$, $f$ running over all Weyl fields and $Y_u,Y_d,Y_e,Y_\nu$ the $3\times3$ Yukawa matrices. With $\langle\Phi\rangle=(0,v/\sqrt2)$ and the inputs $G_F=1.1663788\times10^{-5}$ GeV$^{-2}$, $m_W=80.3692$ GeV, $m_Z=91.1880$ GeV, $m_H=125.20$ GeV (PDG 2025), at tree level $v=(\sqrt2G_F)^{-1/2}=246.220$ GeV, $g=2m_W/v=0.65283$, $g'=0.34995$, $\sin^2\theta_W=1-m_W^2/m_Z^2=0.22321$, $e=0.30843$ and $\lambda=m_H^2/(2v^2)=0.12928$ (section 17.3).

\Needspace{12\baselineskip}\section*{9. dirac16complex at the brane: four brane Dirac spinors}\addcontentsline{toc}{section}{9. dirac16complex at the brane: four brane Dirac spinors}
\Needspace{8\baselineskip}\subsection*{9.1 The decomposition}\addcontentsline{toc}{subsection}{9.1 The decomposition}
dirac16complex is a 16-component field with real $16\times16$ Dirac matrices $\Gamma_A$, $\Gamma_A\Gamma_B+\Gamma_B\Gamma_A=2\eta_{AB}1_{16}$. At the brane the eight directions split into four along $W$ and four transverse to it. With the complex $4\times4$ Dirac matrices of the brane (Dirac representation) and four anti-Hermitian $4\times4$ matrices $\rho_k$,
\begin{align*}
\sigma_1&=\begin{pmatrix}0&1\\1&0\end{pmatrix},\qquad \sigma_2=\begin{pmatrix}0&-i\\i&0\end{pmatrix},\qquad \sigma_3=\begin{pmatrix}1&0\\0&-1\end{pmatrix}\qquad(\text{the Pauli matrices}), \\
\gamma_0&=\mathrm{diag}(1,1,-1,-1),\qquad \gamma_i=\begin{pmatrix}0&\sigma_i\\-\sigma_i&0\end{pmatrix}\ \ (i=1,2,3),\qquad \gamma_5=i\gamma_0\gamma_1\gamma_2\gamma_3, \\
\rho_4&=i\gamma_0,\qquad \rho_5=-\gamma_1,\qquad \rho_6=-\gamma_2,\qquad \rho_7=-\gamma_3, \\
R_\mu&=\gamma_\mu\otimes1_4\ \ (\mu=0,\dots,3),\qquad R_k=\gamma_5\otimes\rho_k\ \ (k=4,\dots,7),\qquad R_AR_B+R_BR_A=2\eta_{AB}1_{16} .
\end{align*}
Both sets of sixteen-dimensional matrices are irreducible representations of the same Clifford algebra, so (Schur’s lemma) there is a unitary matrix $U$, unique up to a phase, with $\Gamma_AU=UR_A$ for every $A$; 7+1.nb obtains it as the Schur average over the 256 Clifford products,
\begin{equation*}
U\propto\sum_{K}\Gamma^{(K)}\,s\,\big(R^{(K)}\big)^{-1},\qquad \Gamma^{(K)}=\Gamma_{A_1}\cdots\Gamma_{A_r},\quad R^{(K)}=R_{A_1}\cdots R_{A_r}\quad(A_1<\dots<A_r),
\end{equation*}
with $s$ the first matrix unit for which the sum is non-zero, normalized to $U^\dagger U=1_{16}$. Writing
\begin{equation*}
\psi=U\zeta,\qquad \zeta=\sum_{a=1}^4\psi_a\otimes e_a\in\mathbb C^4\otimes\mathbb C^4,
\end{equation*}
dirac16complex at the brane is FOUR brane Dirac spinors $\psi_a$ ($a$ the transverse spinor index, $e_a$ the unit vectors of $\mathbb C^4$). With $\bar\psi=\psi^\dagger\Gamma_0$, $\bar\psi_a=\psi_a^\dagger\gamma_0$, $\Gamma^\mu=\eta^{\mu\mu}\Gamma_\mu$ and $\gamma^\mu=\eta^{\mu\mu}\gamma_\mu$:
\begin{equation*}
\bar\psi\psi'=\sum_{a=1}^4\bar\psi_a\psi'_a,\qquad \bar\psi\,\Gamma^\mu\psi'=\sum_{a=1}^4\bar\psi_a\gamma^\mu\psi'_a\quad(\mu=0,\dots,3).
\end{equation*}
\Needspace{6\baselineskip}\subsubsection*{Example 9.1 — the explicit matrices}
Every $\Gamma_A$ of 7+1.nb is a signed permutation matrix: column $j$ has a single non-zero entry, $\mathrm{sign}(c)$ in row $|c|$, with the codes $c$ listed below (rows $A=0,\dots,7$, columns $j=1,\dots,16$):

{\footnotesize
\begin{longtable}{>{\raggedright\arraybackslash}p{\dimexpr 0.0600\linewidth-2\tabcolsep\relax}>{\raggedright\arraybackslash}p{\dimexpr 0.0587\linewidth-2\tabcolsep\relax}>{\raggedright\arraybackslash}p{\dimexpr 0.0587\linewidth-2\tabcolsep\relax}>{\raggedright\arraybackslash}p{\dimexpr 0.0587\linewidth-2\tabcolsep\relax}>{\raggedright\arraybackslash}p{\dimexpr 0.0587\linewidth-2\tabcolsep\relax}>{\raggedright\arraybackslash}p{\dimexpr 0.0587\linewidth-2\tabcolsep\relax}>{\raggedright\arraybackslash}p{\dimexpr 0.0587\linewidth-2\tabcolsep\relax}>{\raggedright\arraybackslash}p{\dimexpr 0.0587\linewidth-2\tabcolsep\relax}>{\raggedright\arraybackslash}p{\dimexpr 0.0587\linewidth-2\tabcolsep\relax}>{\raggedright\arraybackslash}p{\dimexpr 0.0587\linewidth-2\tabcolsep\relax}>{\raggedright\arraybackslash}p{\dimexpr 0.0587\linewidth-2\tabcolsep\relax}>{\raggedright\arraybackslash}p{\dimexpr 0.0587\linewidth-2\tabcolsep\relax}>{\raggedright\arraybackslash}p{\dimexpr 0.0587\linewidth-2\tabcolsep\relax}>{\raggedright\arraybackslash}p{\dimexpr 0.0587\linewidth-2\tabcolsep\relax}>{\raggedright\arraybackslash}p{\dimexpr 0.0587\linewidth-2\tabcolsep\relax}>{\raggedright\arraybackslash}p{\dimexpr 0.0587\linewidth-2\tabcolsep\relax}>{\raggedright\arraybackslash}p{\dimexpr 0.0587\linewidth-2\tabcolsep\relax}}
\toprule
\textbf{$A$} & \textbf{1} & \textbf{2} & \textbf{3} & \textbf{4} & \textbf{5} & \textbf{6} & \textbf{7} & \textbf{8} & \textbf{9} & \textbf{10} & \textbf{11} & \textbf{12} & \textbf{13} & \textbf{14} & \textbf{15} & \textbf{16} \\
\midrule\endfirsthead
\toprule
\textbf{$A$} & \textbf{1} & \textbf{2} & \textbf{3} & \textbf{4} & \textbf{5} & \textbf{6} & \textbf{7} & \textbf{8} & \textbf{9} & \textbf{10} & \textbf{11} & \textbf{12} & \textbf{13} & \textbf{14} & \textbf{15} & \textbf{16} \\
\midrule\endhead
\endfoot
\bottomrule\endlastfoot
0 & 9 & 10 & 11 & 12 & 13 & 14 & 15 & 16 & 1 & 2 & 3 & 4 & 5 & 6 & 7 & 8 \\
1 & -8 & -7 & 6 & 5 & -4 & -3 & 2 & 1 & 16 & 15 & -14 & -13 & 12 & 11 & -10 & -9 \\
2 & 7 & -8 & -5 & 6 & 3 & -4 & -1 & 2 & -15 & 16 & 13 & -14 & -11 & 12 & 9 & -10 \\
3 & -6 & 5 & -8 & 7 & -2 & 1 & -4 & 3 & 14 & -13 & 16 & -15 & 10 & -9 & 12 & -11 \\
4 & -2 & 1 & 4 & -3 & 6 & -5 & -8 & 7 & 10 & -9 & -12 & 11 & -14 & 13 & 16 & -15 \\
5 & 3 & 4 & -1 & -2 & -7 & -8 & 5 & 6 & -11 & -12 & 9 & 10 & 15 & 16 & -13 & -14 \\
6 & 4 & -3 & 2 & -1 & -8 & 7 & -6 & 5 & -12 & 11 & -10 & 9 & 16 & -15 & 14 & -13 \\
7 & 5 & 6 & 7 & 8 & -1 & -2 & -3 & -4 & -13 & -14 & -15 & -16 & 9 & 10 & 11 & 12 \\*
\end{longtable}}
(for example $\Gamma_0$ swaps the components 1–8 with 9–16). The intertwiner has a simple form: $2\sqrt2\,U$ has only the entries $0,\pm1,\pm i$, eight non-zero entries in each row and column:
\begin{equation*}
2\sqrt2\,U=\left(\begin{smallmatrix}0 & 1 & 0 & -i & -i & 0 & 1 & 0 & 0 & 1 & 0 & i & i & 0 & 1 & 0 \\ 0 & i & 0 & 1 & -1 & 0 & -i & 0 & 0 & i & 0 & -1 & 1 & 0 & -i & 0 \\ -i & 0 & 1 & 0 & 0 & -1 & 0 & i & i & 0 & 1 & 0 & 0 & -1 & 0 & -i \\ 1 & 0 & i & 0 & 0 & i & 0 & 1 & -1 & 0 & i & 0 & 0 & i & 0 & -1 \\ 0 & i & 0 & -1 & -1 & 0 & i & 0 & 0 & -i & 0 & -1 & -1 & 0 & -i & 0 \\ 0 & -1 & 0 & -i & i & 0 & 1 & 0 & 0 & 1 & 0 & -i & i & 0 & -1 & 0 \\ -1 & 0 & i & 0 & 0 & -i & 0 & 1 & -1 & 0 & -i & 0 & 0 & i & 0 & 1 \\ -i & 0 & -1 & 0 & 0 & -1 & 0 & -i & -i & 0 & 1 & 0 & 0 & 1 & 0 & -i \\ 0 & 1 & 0 & -i & -i & 0 & 1 & 0 & 0 & -1 & 0 & -i & -i & 0 & -1 & 0 \\ 0 & i & 0 & 1 & -1 & 0 & -i & 0 & 0 & -i & 0 & 1 & -1 & 0 & i & 0 \\ -i & 0 & 1 & 0 & 0 & -1 & 0 & i & -i & 0 & -1 & 0 & 0 & 1 & 0 & i \\ 1 & 0 & i & 0 & 0 & i & 0 & 1 & 1 & 0 & -i & 0 & 0 & -i & 0 & 1 \\ 0 & i & 0 & -1 & -1 & 0 & i & 0 & 0 & i & 0 & 1 & 1 & 0 & i & 0 \\ 0 & -1 & 0 & -i & i & 0 & 1 & 0 & 0 & -1 & 0 & i & -i & 0 & 1 & 0 \\ -1 & 0 & i & 0 & 0 & -i & 0 & 1 & 1 & 0 & i & 0 & 0 & -i & 0 & -1 \\ -i & 0 & -1 & 0 & 0 & -1 & 0 & -i & i & 0 & -1 & 0 & 0 & -1 & 0 & i\end{smallmatrix}\right)
\end{equation*}
(both printed by section 25.7 of 7+1.nb and recomputed by the Python program, which also checks with the null space of $X\mapsto(\Gamma_AX-XR_A)_A$ that every intertwiner is a multiple of $U$).

\Needspace{6\baselineskip}\subsubsection*{Example 9.2 — one brane spinor}
Take the brane spinor $u=(1,2,0,-1)$ in the first transverse slot, $\zeta=u\otimes e_1$, $\psi=U\zeta$. Then
\begin{equation*}
\bar\psi\psi=\bar uu=|u_1|^2+|u_2|^2-|u_3|^2-|u_4|^2=4,\qquad \big(\bar\psi\Gamma^\mu\psi\big)_{\mu=0,\dots,3}=\big(\bar u\gamma^\mu u\big)_{\mu=0,\dots,3}=(6, 2, 0, -4)
\end{equation*}
— the sixteen-component bilinears reduce to those of one ordinary Dirac spinor (section 25.7).

\Needspace{8\baselineskip}\subsection*{9.2 Two kinds of rotations act on two different indices}\addcontentsline{toc}{subsection}{9.2 Two kinds of rotations act on two different indices}
Products of two Dirac matrices generate rotations. From $\Gamma_AU=UR_A$:
\begin{equation*}
\Gamma_\mu\Gamma_\nu=U\,(\gamma_\mu\gamma_\nu\otimes1_4)\,U^\dagger,\qquad \Gamma_k\Gamma_l=U\,(1_4\otimes\rho_k\rho_l)\,U^\dagger,\qquad \Gamma_\mu\Gamma_k=U\,(\gamma_\mu\gamma_5\otimes\rho_k)\,U^\dagger,
\end{equation*}
(using $\gamma_5^2=1_4$). The Lorentz group of the brane acts on the brane spinor index only, the rotations about the brane on the transverse index $a$ only, and the generators $\Gamma_\mu\Gamma_k$, which would rotate the brane into the transverse directions, act on both — they are the symmetries the brane breaks (section 25.7 checks all 48 products with $\mu,\nu=0,\dots,3$ and $k,l=4,\dots,7$).

\Needspace{8\baselineskip}\subsection*{9.3 The transverse index carries \texorpdfstring{$Spin(4)=SU(2)\times SU(2)$}{Spin(4)=SU(2) SU(2)}}\addcontentsline{toc}{subsection}{9.3 The transverse index carries \texorpdfstring{$Spin(4)=SU(2)\times SU(2)$}{Spin(4)=SU(2) SU(2)}}
The product $\rho_{4567}=\rho_4\rho_5\rho_6\rho_7$ equals $-\gamma_5$ (as a $4\times4$ matrix), is Hermitian with the eigenvalues $+1,+1,-1,-1$, anticommutes with every $\rho_k$ and commutes with every $\rho_k\rho_l$. For each even permutation $(k,l,m,n)$ of $(4,5,6,7)$, $\rho_m\rho_n=-\rho_k\rho_l\rho_{4567}$, so on the eigenspace $\rho_{4567}=s$
\begin{equation*}
\big(\rho_k\rho_l+s\,\rho_m\rho_n\big)P_s=0,\qquad P_s=\tfrac12(1_4+s\,\rho_{4567}),
\end{equation*}
while $\rho_k\rho_l-s\,\rho_m\rho_n$ acts there with rank 2. The six rotations about the brane combine into two commuting sets of three (the self-dual and anti-self-dual combinations), each an $SU(2)$, and each eigenspace of $\rho_{4567}$ feels only one of them: the four values of $a$ form the representation $(2,1)\oplus(1,2)$ of $Spin(4)$ (section 25.7).

\Needspace{8\baselineskip}\subsection*{9.4 Chirality at the brane}\addcontentsline{toc}{subsection}{9.4 Chirality at the brane}
The eight-dimensional chirality is $\chi=i\Gamma_0\Gamma_1\cdots\Gamma_7$ (the factor $i$ is needed because $(\Gamma_0\cdots\Gamma_7)^2=-1_{16}$: in 7+1 dimensions there is no real chirality), with $\chi^2=1_{16}$ and the projector $P_L=\frac12(1-\chi)$ onto the left-handed half. From $\Gamma_AU=UR_A$, $R_0R_1R_2R_3=-i\gamma_5\otimes1_4$ and $R_4R_5R_6R_7=1_4\otimes\rho_{4567}$:
\begin{equation*}
\chi=U(\gamma_5\otimes\rho_{4567})U^\dagger,\qquad P_L=U\Big(\tfrac{1-\gamma_5}2\otimes P_+ +\tfrac{1+\gamma_5}2\otimes P_-\Big)U^\dagger :
\end{equation*}
the left-handed half $\psi_L$ (eight components) consists of two brane-LEFT-handed Weyl spinors ($\gamma_5=-1$) with $\rho_{4567}=+1$ and two brane-RIGHT-handed ones ($\gamma_5=+1$) with $\rho_{4567}=-1$; each of the two parts has rank 4 (two Weyl spinors of two components each; section 25.7). Eight-dimensional chirality is not four-dimensional chirality: it is the product of the brane chirality and the transverse chirality.

\Needspace{8\baselineskip}\subsection*{9.5 The Kaluza–Klein tower at the brane}\addcontentsline{toc}{subsection}{9.5 The Kaluza–Klein tower at the brane}
A mode $\psi\propto e^{i\vec q\cdot\vec y}$, $\vec q=2\pi\vec n/L_c$ ($\vec n\in\mathbb Z^4$), of the static torus satisfies at the brane
\begin{equation*}
\big(\gamma\cdot p\otimes1_4+X-m\big)\zeta=0,\qquad X=\gamma_5\otimes\frac{\vec\rho\cdot\vec q}{b_0},\qquad X^2=-\frac{|\vec q|^2}{b_0^2}1_{16},\qquad m_{\vec n}=\sqrt{m^2+\frac{4\pi^2|\vec n|^2}{\ell_\perp^2}},
\end{equation*}
($\gamma\cdot p=p_0\gamma_0-\sum_ip_i\gamma_i$; $X$ anticommutes with every $\gamma_\mu\otimes1_4$): each momentum vector $\vec n$ gives four brane Dirac fermions of mass $m_{\vec n}$. Every mode has the same value at the brane, so the bulk field seen from the brane is the whole tower with equal weights,
\begin{equation*}
\psi(x,\vec y=0)=\frac1{\sqrt{V_\perp}}\sum_{\vec n\in\mathbb Z^4}\chi_{\vec n}(x),
\end{equation*}
($\chi_{\vec n}$ the four-dimensional fields, four Dirac spinors each). The lowest levels for $m=0$:

{\footnotesize
\begin{longtable}{>{\raggedright\arraybackslash}p{\dimexpr 0.1600\linewidth-2\tabcolsep\relax}>{\raggedright\arraybackslash}p{\dimexpr 0.2600\linewidth-2\tabcolsep\relax}>{\raggedright\arraybackslash}p{\dimexpr 0.2800\linewidth-2\tabcolsep\relax}>{\raggedright\arraybackslash}p{\dimexpr 0.3000\linewidth-2\tabcolsep\relax}}
\toprule
\textbf{$k=|\vec n|^2$} & \textbf{momentum vectors $r_4(k)$} & \textbf{brane Dirac fermions $4r_4(k)$} & \textbf{mass $2\pi\sqrt k/\ell_\perp$ (MeV)} \\
\midrule\endfirsthead
\toprule
\textbf{$k=|\vec n|^2$} & \textbf{momentum vectors $r_4(k)$} & \textbf{brane Dirac fermions $4r_4(k)$} & \textbf{mass $2\pi\sqrt k/\ell_\perp$ (MeV)} \\
\midrule\endhead
\endfoot
\bottomrule\endlastfoot
1 & 8 & 32 & 4.02652 \\
2 & 24 & 96 & 5.69437 \\
3 & 32 & 128 & 6.97415 \\
4 & 24 & 96 & 8.05305 \\
5 & 48 & 192 & 9.00358 \\
6 & 96 & 384 & 9.86293 \\
7 & 64 & 256 & 10.6532 \\
8 & 24 & 96 & 11.3887 \\*
\end{longtable}}
($r_4(k)=8\sigma(k)-32\sigma(k/4)$ is the number of $\vec n\in\mathbb Z^4$ with $|\vec n|^2=k$, $\sigma$ the divisor sum.)

\Needspace{12\baselineskip}\section*{10. The portals: where the bulk meets the brane}\addcontentsline{toc}{section}{10. The portals: where the bulk meets the brane}
The couplings between the Standard Model and dirac16complex (an invented interaction of the model; separate couplings for the two eight-dimensional chiral halves) are local at the brane, with $\psi_a(x)$ the brane spinors of $\psi(x,\vec y=0)$:

\begin{itemize}
\item \textbf{Neutrino (Yukawa) portal:} $-\sum_{i,s}\big[y^L_{is}\bar L_i\tilde\Phi P_R\psi^{(L)}_s+y^R_{is}\bar L_i\tilde\Phi P_R\psi^{(R)}_s\big]+\text{h.c.}$, with $\psi^{(L)}_s$ ($\psi^{(R)}_s$), $s=1,2$, the brane-right-handed parts of the left (right) half of dirac16complex and $P_R=\frac12(1+\gamma_5)$. After symmetry breaking it is a Dirac mass $\mu_{is}=y_{is}v/(\sqrt2\sqrt{V_\perp})$ between $\nu_{L,i}$ and EVERY Kaluza–Klein level.
\item \textbf{Vector (current) portal:} $-G^f_P(\bar\psi\Gamma^\mu\psi)|_W(\bar f\gamma_\mu f)$ for Standard Model fermions $f$.
\item \textbf{Higgs (scalar) portal:} $-\lambda_P\Phi^\dagger\Phi(\bar\psi\psi)|_W$.
\end{itemize}
Because the brane sits at one point of the torus, these couplings do not conserve Kaluza–Klein momentum: the product of three modes at the brane is $(\phi_{\vec n}\phi_{\vec m}\phi^*_{\vec l})(0)=V_\perp^{-3/2}$ for ALL $\vec n,\vec m,\vec l$ (with $\phi_{\vec n}=e^{i\vec q_{\vec n}\cdot\vec y}/\sqrt{V_\perp}$), while in the bulk the overlap $\int b^4d^4y\,\phi_{\vec n}\phi_{\vec m}\phi^*_{\vec l}=\delta_{\vec n+\vec m,\vec l}/\sqrt{V_\perp}$ conserves it (section 25.5). The Higgs portal, for example, couples every pair of levels, $\bar\psi\psi|_W=\sum_{\vec n,\vec n'}\bar\chi_{\vec n}\chi_{\vec n'}/V_\perp$.

\Needspace{6\baselineskip}\subsubsection*{Example 10.1 — natural couplings are diluted by the volume}
For couplings of natural size at $M_8$ — $\lambda_P=c_P/M_8^5$, $y=c_y/M_8^2$, $G_P=c_V/M_8^6$ — the four-dimensional couplings of a level are
\begin{equation*}
\lambda_{\rm eff}=\frac{\lambda_P}{V_\perp}=\frac{c_PM_8}{M_{\rm Pl}^2},\qquad y_{\rm eff}=\frac y{\sqrt{V_\perp}}=\frac{c_yM_8}{M_{\rm Pl}},\qquad G_{\rm eff}=\frac{G_P}{V_\perp}=\frac{c_V}{M_{\rm Pl}^2}\qquad\Big(V_\perp=\frac{M_{\rm Pl}^2}{M_8^6}\Big),
\end{equation*}
i.e. $\lambda_{\rm eff}=1.69\times10^{-33}\,c_P$ GeV$^{-1}$ and $y_{\rm eff}=4.11\times10^{-15}\,c_y$, and the portal Dirac mass $\mu_P=y_{\rm eff}v/\sqrt2=7.15\times10^{-13}\,c_y$ GeV (section 24.3 of 7+1.nb): the brane talks to each level very weakly, but to very many levels.

\Needspace{6\baselineskip}\subsubsection*{Example 10.2 — a neutrino mass matrix with a portal}
In the basis $(\nu_L,\psi_L)\times(\nu_R,\psi_R)$ the Yukawa portal gives the zero mode the Dirac mass matrix
\begin{equation*}
\mathcal M=\begin{pmatrix}m_D&\mu\\0&m\,1_4\end{pmatrix},\qquad m_D=U_{\rm PMNS}\,\mathrm{diag}(m_1,m_2,m_3),\qquad 1-NN^\dagger\simeq\frac{\mu\mu^\dagger}{m^2},
\end{equation*}
($N$ the light $3\times3$ block of the mixing matrix, $U_{\rm PMNS}$ the lepton mixing matrix and $m_1,m_2,m_3$ the light neutrino masses of section 17 of 7+1.nb). Section 18.2 of 7+1.nb takes $m=1$ GeV and
\begin{equation*}
\mu=10^4\ \text{eV}\times\begin{pmatrix}3&1&0&0\\1&2&0&0\\-2&1&0&0\end{pmatrix}
\end{equation*}
and finds, from the singular values of $\mathcal M$,
\begin{align*}
m_{\rm light}&=(0,\ 0.008654,\ 0.05034)\ \text{eV},\qquad (1-NN^\dagger)_{ee}=1.000\times10^{-9}, \\
m_{\rm heavy}&=(1.00000000000,\ 1.00000000000,\ 1.00000000025,\ 1.00000000075)\ \text{GeV}:
\end{align*}
the heavy states are mostly dirac16complex, and the light mixing matrix is no longer unitary, by $(\mu\mu^\dagger)_{ee}/m^2=10^{-9}$ in the $ee$ entry.

\Needspace{12\baselineskip}\section*{11. The brane in the three regimes of the model}\addcontentsline{toc}{section}{11. The brane in the three regimes of the model}
\textbf{Primordial regime.} Near the singularity the bulk is dominated by Lovelock gravity and a gas of dirac16complex; even at the temperature $T=M_8$ and the expansion rate $H=M_8$ the Standard Model plasma would be the fraction $(\pi^2g_*/30)M_8^4/(3M_{\rm Pl}^2M_8^2)=2.0\times10^{-28}$ of the energy ($g_*=106.75$; section 21.4): the brane is gravitationally negligible.

\textbf{The normalcy temperature.} A hot brane loses energy by emitting Kaluza–Klein modes into the bulk. The relative rate is about $T^7/M_8^6$ — the number $V_\perp T^4$ of levels lighter than $T$ times the rate $T^3/M_{\rm Pl}^2$ per level — and it equals the expansion rate $T^2/M_{\rm Pl}$ at
\begin{equation*}
T_*=\Big(\frac{M_8^6}{M_{\rm Pl}}\Big)^{1/5}=13.26\ \text{GeV}
\end{equation*}
(ESTIMATE, numerical factors omitted; section 22.2). Below $T_*$ the brane plasma cools by the expansion; the modern regime of 7+1.nb starts there.

\textbf{Modern regime.} With the torus static the brane is an ordinary Friedmann universe. The radiation era uses the relativistic degrees of freedom of the Standard Model, $\rho_4=\frac{\pi^2}{30}g_*T^4$, $s=\frac{2\pi^2}{45}g_{*s}T^3$, $H_a=(\pi^2g_*/90)^{1/2}T^2/M_{\rm Pl}$ (section 22.3):

{\footnotesize
\begin{longtable}{>{\raggedright\arraybackslash}p{\dimexpr 0.3000\linewidth-2\tabcolsep\relax}>{\raggedright\arraybackslash}p{\dimexpr 0.3500\linewidth-2\tabcolsep\relax}>{\raggedright\arraybackslash}p{\dimexpr 0.3500\linewidth-2\tabcolsep\relax}}
\toprule
\textbf{$T$} & \textbf{$g_*(T)$} & \textbf{$g_{*s}(T)$} \\
\midrule\endfirsthead
\toprule
\textbf{$T$} & \textbf{$g_*(T)$} & \textbf{$g_{*s}(T)$} \\
\midrule\endhead
\endfoot
\bottomrule\endlastfoot
1 TeV & 106.79 & 106.78 \\
10 GeV & 86.296 & 86.205 \\
1 GeV & 76.418 & 75.577 \\
300 MeV & 64.646 & 64.217 \\
100 MeV & 18.315 & 17.950 \\
10 MeV & 10.837 & 10.835 \\
1 MeV & 10.634 & 10.636 \\
100 keV & 4.3270 & 4.8144 \\
1 keV & 3.3830 & 3.9375 \\*
\end{longtable}}
and the time–temperature relation gives $t=0.743$ s at $T=1$ MeV and $13162$ s at 10 keV. Today the brane universe is 13.791 Gyr old (section 22.4: $\Omega_r=9.210\times10^{-5}$, $\Omega_\Lambda=0.68491$, matter–radiation equality at $z=3419.4$), and the radion is displaced by only $1.7\times10^{-86}$ (section 22.4).

\textbf{Far future.} $\Lambda_4$ dominates, the brane expands like de Sitter space with $H_\infty=H_0\Omega_\Lambda^{1/2}=1.808\times10^{-18}$ s$^{-1}$ ($a\propto\sinh^{2/3}(\frac32H_\infty t)$ for matter $+$ $\Lambda$), and the brane plasma dilutes like $e^{-4H_\infty t}$ (section 23).

\Needspace{12\baselineskip}\section*{12. The brane on a 4-sphere instead of the 4-torus}\addcontentsline{toc}{section}{12. The brane on a 4-sphere instead of the 4-torus}
7+1.nb uses a torus. To see which properties of the brane depend on that choice, put the brane into a 4-SPHERE of radius $r$ (section 26 of 7+1.nb): the extra coordinates are the angles $x_4,x_5,x_6\in[0,\pi]$, $x_7\in[0,2\pi)$, the metric of the extra space is $r^2\hat\gamma$ with
\begin{equation*}
\hat\gamma=\mathrm{diag}\big(1,\ \sin^2x_4,\ \sin^2x_4\sin^2x_5,\ \sin^2x_4\sin^2x_5\sin^2x_6\big),\qquad V=\frac{8\pi^2}3r^4,\qquad R_{ij}=3\hat\gamma_{ij},
\end{equation*}
and the brane sits at the NORTH POLE $x_4=0$, the point $(0,0,0,0,r)$ of the embedding of the sphere in $\mathbb R^5$
\begin{equation*}
(X_1,\dots,X_5)=r\big(\sin x_4\sin x_5\sin x_6\cos x_7,\ \sin x_4\sin x_5\sin x_6\sin x_7,\ \sin x_4\sin x_5\cos x_6,\ \sin x_4\cos x_5,\ \cos x_4\big).
\end{equation*}
\Needspace{8\baselineskip}\subsection*{12.1 What stays the same: the induced metric and the radion coupling}\addcontentsline{toc}{subsection}{12.1 What stays the same: the induced metric and the radion coupling}
With the radion $\beta$ ($r=r_0e^\beta$) the bulk metric is $ds^2=e^{-4\beta}g_E-e^{2\beta}r_0^2\hat\gamma$, and at the north pole the induced metric is
\begin{equation*}
h_{\mu\nu}=e^{-4\beta}g_{E\mu\nu}\qquad\text{— the same as at }\vec y=0\text{ of the torus.}
\end{equation*}
So the coupling $\phi\,T^\mu{}_\mu/(\sqrt3M_{\rm Pl})$ of section 7.3, the fifth-force strength $\frac23$ of a massless radion and Newton’s constant $8\pi G_4=\kappa/V$ are the same on both spaces; only the volume formula changes. With the volume of the torus today, $V=V_\perp$, the sphere has the radius $r_{S^4}=(3V_\perp/(8\pi^2))^{1/4}=1.3595\times10^{-13}$ m.

\Needspace{8\baselineskip}\subsection*{12.2 What changes: symmetry}\addcontentsline{toc}{subsection}{12.2 What changes: symmetry}
The isometries of $S^4$ are the 10 rotations $L_{ab}=X_a\partial_b-X_b\partial_a$ of $\mathbb R^5$ ($SO(5)$). At the north pole the 6 with $a,b\le4$ vanish — the brane keeps $SO(4)$, which acts on the brane fields as rotations about the brane — and the 4 with $b=5$ move it (section 26.4 counts them). On the torus the brane keeps the 384 signed permutations and breaks the 4 translations. In both cases the broken symmetries mean that Kaluza–Klein quantum numbers (momentum $\vec n$ on $T^4$, angular momentum $l$ on $S^4$) are not conserved at the brane.

\Needspace{8\baselineskip}\subsection*{12.3 What changes: the Kaluza–Klein states the brane sees}\addcontentsline{toc}{subsection}{12.3 What changes: the Kaluza–Klein states the brane sees}
A bulk scalar on $S^4$ has the levels $\mu_l=\sqrt{l(l+3)}/r$ with the multiplicity $d_l=(2l+3)(l+1)(l+2)/6$ (the spherical harmonics of degree $l$). At the north pole only one harmonic of each level is non-zero, the ZONAL harmonic
\begin{equation*}
Z_l(x_4)=\sqrt{\frac{d_l}V}\,\frac{C_l^{(3/2)}(\cos x_4)}{C_l^{(3/2)}(1)},\qquad C_l^{(3/2)}(1)=\frac{(l+1)(l+2)}2,\qquad Z_l(0)=\sqrt{\frac{d_l}V},
\end{equation*}
($C_l^{(3/2)}$ the Gegenbauer polynomial; normalized, $\int_{S^4}Z_l^2=1$). A brane field coupled to the bulk field at the brane therefore couples to ONE state per level with the strength $\sqrt{d_l/V}$; on the torus it couples to the combination $\sum_{|\vec n|^2=k}\psi_{\vec n}$ of each shell with the strength $\sqrt{r_4(k)/V}$. \textbf{Example 12.3:} the first levels with $r=r_{S^4}$:

{\footnotesize
\begin{longtable}{>{\raggedright\arraybackslash}p{\dimexpr 0.1500\linewidth-2\tabcolsep\relax}>{\raggedright\arraybackslash}p{\dimexpr 0.1500\linewidth-2\tabcolsep\relax}>{\raggedright\arraybackslash}p{\dimexpr 0.3500\linewidth-2\tabcolsep\relax}>{\raggedright\arraybackslash}p{\dimexpr 0.3500\linewidth-2\tabcolsep\relax}}
\toprule
\textbf{$l$} & \textbf{$d_l$} & \textbf{$\mu_l=\sqrt{l(l+3)}/r$ (MeV)} & \textbf{brane coupling} \\
\midrule\endfirsthead
\toprule
\textbf{$l$} & \textbf{$d_l$} & \textbf{$\mu_l=\sqrt{l(l+3)}/r$ (MeV)} & \textbf{brane coupling} \\
\midrule\endhead
\endfoot
\bottomrule\endlastfoot
$0$ & 1 & 0 & $\sqrt{1/V}$ \\
$1$ & 5 & 2.9030 & $\sqrt{5/V}$ \\
$2$ & 14 & 4.5901 & $\sqrt{14/V}$ \\
$3$ & 30 & 6.1582 & $\sqrt{30/V}$ \\*
\end{longtable}}
while the torus of the same volume has its first level at $2\pi/\ell_\perp=4.0265$ MeV with $r_4(1)=8$ states.

\Needspace{8\baselineskip}\subsection*{12.4 What changes: dirac16complex at the brane}\addcontentsline{toc}{subsection}{12.4 What changes: dirac16complex at the brane}
On the torus the Dirac operator of the extra space has zero modes (the constant spinors): for $m=0$ dirac16complex gives four MASSLESS brane Dirac spinors (section 9). On $S^4$ it has none: its eigenvalues are $\pm i(l+2)/r$ ($l=0,1,\dots$, degeneracy $4\binom{l+3}3$ for each sign), and the lowest, $\pm2i/r$, belong to the Killing spinors ($\tilde\gamma_a=i\rho_{a+3}$, frame $e^a=(r,r\sin x_4,\dots)dx^{a+3}$)
\begin{equation*}
\varepsilon_s=e^{s\frac i2x_4\tilde\gamma_1}e^{\frac12x_5\tilde\gamma_1\tilde\gamma_2}e^{\frac12x_6\tilde\gamma_2\tilde\gamma_3}e^{\frac12x_7\tilde\gamma_3\tilde\gamma_4}\varepsilon_0,\qquad \nabla_i\varepsilon_s=s\frac i{2r}e_i{}^a\tilde\gamma_a\varepsilon_s\qquad(s=\pm1),
\end{equation*}
which saturate Friedrich’s bound $|\lambda|^2\ge\frac{nR_{\rm int}}{4(n-1)}=\frac4{r^2}$. The lightest dirac16complex states seen at a sphere brane are eight Dirac fermions of mass $\sqrt{m^2+4/r^2}\ge2.9030$ MeV: a sphere brane gets no massless fermions from the bulk field.

\Needspace{8\baselineskip}\subsection*{12.5 What changes: holding the size}\addcontentsline{toc}{subsection}{12.5 What changes: holding the size}
The sphere is curved, so its size costs energy: the radion potential contains $-6M_{\rm Pl}^2e^{-6\beta}/r_0^2$. A flat four-dimensional vacuum needs, with the stabilizing energy density $U=U_0+U_1\beta+\frac{K_U}2\beta^2$ and the bulk cosmological constant $\Lambda$,
\begin{equation*}
(\Lambda+\kappa U_0)\,r_0^2=6,\qquad \kappa U_1\,r_0^2=-12,\qquad m_r^2=\frac{\kappa K_U}{12}-2\,\frac1{r_0^2}\qquad\big(T^4:\ \Lambda+\kappa U_0=0,\ U_1=0,\ m_r^2=\tfrac{\kappa K_U}{12}\big),
\end{equation*}
and a sphere held by $\Lambda$ alone ($dS_4\times S^4$, $\Lambda r_0^2=9$) sits at a maximum, $m^2=-6\,H^2$ (section 26.2). The brane itself does not change these conditions as long as its tension is small (its energy smeared over the sphere adds $f_W/V$).

\Needspace{8\baselineskip}\subsection*{12.6 Side by side}\addcontentsline{toc}{subsection}{12.6 Side by side}
{\footnotesize
\begin{longtable}{>{\raggedright\arraybackslash}p{\dimexpr 0.3000\linewidth-2\tabcolsep\relax}>{\raggedright\arraybackslash}p{\dimexpr 0.3500\linewidth-2\tabcolsep\relax}>{\raggedright\arraybackslash}p{\dimexpr 0.3500\linewidth-2\tabcolsep\relax}}
\toprule
\textbf{} & \textbf{brane in $T^4$ (7+1.nb)} & \textbf{brane at the north pole of $S^4$} \\
\midrule\endfirsthead
\toprule
\textbf{} & \textbf{brane in $T^4$ (7+1.nb)} & \textbf{brane at the north pole of $S^4$} \\
\midrule\endhead
\endfoot
\bottomrule\endlastfoot
position & $\vec y=0$ & $x_4=0$ \\
induced metric & $e^{-4\beta}g_E$ & $e^{-4\beta}g_E$ \\
radion coupling & $\phi T^\mu{}_\mu/(\sqrt3M_{\rm Pl})$ & the same \\
symmetry kept / broken & 384 signed permutations / 4 translations & $SO(4)$ (6) / 4 of the 10 rotations \\
states per level that couple & 1 combination of $r_4(k)$ & 1 zonal harmonic of $d_l$ \\
coupling strength & $\sqrt{r_4(k)/V}$ & $\sqrt{d_l/V}$ \\
massless dirac16complex states ($m=0$) & 4 brane Dirac spinors & none (lightest $2/r$) \\
size held by & $U$ with $\Lambda+\kappa U_0=0$ & $U$ with $(\Lambda+\kappa U_0)r_0^2=6$, or a flux \\*
\end{longtable}}
\Needspace{12\baselineskip}\section*{13. Flat versus warped branes: Randall–Sundrum}\addcontentsline{toc}{section}{13. Flat versus warped branes: Randall–Sundrum}
The bulk of 7+1 is a PRODUCT: the brane metric does not depend on where along the extra space it is taken. In the warped model of Randall and Sundrum (RS1, five dimensions) it does: $ds^2=e^{-2A(y)}\eta_{\mu\nu}dx^\mu dx^\nu-dy^2$, $y\in[-\pi R,\pi R]$ with $y\simeq-y$, and two branes, at $y=0$ (UV) and $y=\pi R$ (IR, where the Standard Model lives). Section 26.5 of 7+1.nb checks it in the conventions of this document.

\Needspace{8\baselineskip}\subsection*{13.1 Brane tensions and the junction conditions}\addcontentsline{toc}{subsection}{13.1 Brane tensions and the junction conditions}
\begin{equation*}
G^\mu{}_\nu=3A''-6A'^2,\qquad G^y{}_y=-6A'^2;\qquad A=k|y|:\ \Lambda=-6k^2 .
\end{equation*}
The slope $A'$ jumps by $+2k$ at the UV brane and by $-2k$ at the IR brane, so the $\delta$-function part of $G^\mu{}_\nu$ is $3\cdot(\pm2k)\delta(y)$; a brane with tension $f$ ($T^\mu{}_\nu=f\delta(y)\delta^\mu_\nu$) balances it if
\begin{equation*}
f_{\rm UV}=+\frac{6k}\kappa,\qquad f_{\rm IR}=-\frac{6k}\kappa .
\end{equation*}
In terms of the extrinsic curvature $K^\mu{}_\nu=\frac12g^{\mu\alpha}\partial_yg_{\alpha\nu}=-A'\delta^\mu_\nu$ of the brane, the same statement is the Israel junction condition $[K^\mu{}_\nu]-\delta^\mu_\nu[K]=3[A']=\kappa f\,\delta^\mu_\nu$; in the signature $(-,+,+,+,+)$ it reads $[K_{\mu\nu}-g_{\mu\nu}K]=-\kappa S_{\mu\nu}$. The author’s note, which uses $+\kappa S_{\mu\nu}$, obtains a negative UV tension; the check gives the positive UV tension of Randall and Sundrum ($24M^3k$ for the action $2M^3\int\sqrt{|G|}R$, $\kappa=1/(4M^3)$). The brane of 7+1 is totally geodesic ($K=0$, section 4): it needs no such balance, and its tension $f_W$ is an ordinary contribution to the four-dimensional cosmological constant.

\Needspace{8\baselineskip}\subsection*{13.2 The radion of a warped brane}\addcontentsline{toc}{subsection}{13.2 The radion of a warped brane}
With the radius promoted to a field, $ds^2=e^{-2kT(x)|y|}g-T(x)^2dy^2$, the IR brane sees $h=e^{-2k\pi T}g=(\mu/f)^2g$ with the canonical radion $\mu=fe^{-k\pi T}$, $f^2=6/(\kappa k)$ (section 26.5; the note has $12/(\kappa k)$), and $M_{\rm Pl}^2=(1-e^{-2\pi kR})/(\kappa k)$. Its coupling to brane matter is $(\delta\mu/\langle\mu\rangle)T^\mu{}_\mu$ ($\mu=\langle\mu\rangle+\delta\mu$): TeV-strong, for $\langle\mu\rangle=3$ TeV stronger than the flat radion of 7+1 by
\begin{equation*}
\frac{\sqrt3\,M_{\rm Pl}}{\langle\mu\rangle}=1.4058\times10^{15} .
\end{equation*}
A bulk scalar with brane values $v_A$, $v_B$ (Goldberger–Wise) stabilizes it with $m_r^2=\frac43\kappa k^2\epsilon^2v_B^2(\langle\mu\rangle/f)^2$, $\epsilon=m^2/(4k^2)$; the same brane values fix $k\pi\langle T\rangle=(4k^2/m^2)\ln(v_A/v_B)$. The Higgs boson on the IR brane mixes with this radion through $\gamma=v/\langle\mu\rangle\approx0.082$ — the effect that is negligible in 7+1 (section 7.5) is large here.

\Needspace{8\baselineskip}\subsection*{13.3 No warping in eight dimensions with \texorpdfstring{$\Lambda$}{} alone}\addcontentsline{toc}{subsection}{13.3 No warping in eight dimensions with \texorpdfstring{$\Lambda$}{} alone}
Could the 7+1 brane sit in a warped eight-dimensional bulk? For $ds^2=e^{-2A(y)}\eta-dy^2-d\vec z^{\,2}$ (warping along one of the four extra directions) section 26.5 finds
\begin{equation*}
G^\mu{}_\nu=3A''-6A'^2,\qquad G^y{}_y=-6A'^2,\qquad G^z{}_z=4A''-10A'^2,
\end{equation*}
and with only a cosmological constant all eight components must be equal: $A''=0$ and $A'=0$. The five-dimensional equations do not carry over; a warped eight-dimensional brane would need other sources, which 7+1 does not contain.

\Needspace{12\baselineskip}\section*{14. Validity and limitations}\addcontentsline{toc}{section}{14. Validity and limitations}
\begin{itemize}
\item The brane is FIXED at $\vec y=0$: its position is not a dynamical field, so its fluctuations (the “branons” of the literature) are not part of the model; for the metrics of the model this is consistent, because the brane is totally geodesic (section 4).
\item The brane is THIN (a delta function); its own gravitational field, of order one within about $1/M_8$ of it, is not computed (section 4.1); the homogeneous equations use the source smeared over the torus (section 6.1).
\item The brane action contains a tension, the Standard Model and the portals, but no brane-localized Einstein–Hilbert term ($\int_W\sqrt{|h|}R^{(h)}$); gravity on the brane is entirely the bulk gravity seen through $h$ and $V_\perp$. The curvature coupling $\xi R(h)\Phi^\dagger\Phi$ of the Higgs boson (section 7.5) is allowed; it only mixes the Higgs boson with the radion, negligibly.
\item The comparisons with a 4-sphere (section 12) and with warped branes (section 13) are worked out in section 26 of 7+1.nb; the model itself uses the torus and an unwarped bulk.
\item The portals are invented couplings of the model; their sizes are free (natural size: Example 10.1).
\item Numbers assume $M_8=10$ TeV, $M_{\rm Pl}=2.435\times10^{18}$ GeV, $m_r=10$ GeV and the Planck 2018 cosmology.
\end{itemize}
\Needspace{12\baselineskip}\section*{15. Verification}\addcontentsline{toc}{section}{15. Verification}
The worked examples of sections 3–6 and 9 and the estimates of sections 4.1 and 11 are checks of section 25 of 7+1.nb (Wolfram Language 15.0.1, evaluated in a fresh kernel with 0 message cells) and of the independent Python program tools/ref\_torus\_brane.py (sympy, numpy, mpmath). The checks of section 25 (checks 1–14 concern the torus alone, sections 25.1–25.4; checks 15–16 the torus and the brane together, section 25.5; checks 17–27 the world brane, sections 25.6–25.8):

{\footnotesize
\begin{longtable}{>{\raggedright\arraybackslash}p{\dimexpr 0.0500\linewidth-2\tabcolsep\relax}>{\raggedright\arraybackslash}p{\dimexpr 0.7700\linewidth-2\tabcolsep\relax}>{\raggedright\arraybackslash}p{\dimexpr 0.0800\linewidth-2\tabcolsep\relax}>{\raggedright\arraybackslash}p{\dimexpr 0.1000\linewidth-2\tabcolsep\relax}}
\toprule
\textbf{\#} & \textbf{check} & \textbf{result} & \textbf{section} \\
\midrule\endfirsthead
\toprule
\textbf{\#} & \textbf{check} & \textbf{result} & \textbf{section} \\
\midrule\endhead
\endfoot
\bottomrule\endlastfoot
1 & \texttt{auxiliary torus: for the cosmological metric det(-g\_\allowbreak{}kl) == b\textasciicircum{}8 (k, l = 4..7), and the proper volume of the fundamental domain [0, Lc)\textasciicircum{}4, Integrate[Sqrt[det(-g\_\allowbreak{}kl)], \{y, 0, Lc\}\textasciicircum{}4], is b\textasciicircum{}4 Lc\textasciicircum{}4 = V\_\allowbreak{}perp} & True & 25.1 \\
2 & \texttt{only the product b Lc of the auxiliary torus is physical: y = lambda y' maps the torus of period Lc and size b to the one of period Lc/\allowbreak{}lambda and size lambda b; the line element b\textasciicircum{}2 dy.dy, the proper period b Lc and V\_\allowbreak{}perp = b\textasciicircum{}4 Lc\textasciicircum{}4 are unchanged} & True & 25.1 \\
3 & \texttt{symmetries of the auxiliary torus that fix y = 0 (and the brane): the 384 = 2\textasciicircum{}4 4! signed permutation matrices are distinct orthogonal integer matrices (they map Z\textasciicircum{}4 onto itself and preserve |n|\textasciicircum{}2); the shells |n|\textasciicircum{}2 = 1, 2, 3, 4, 9 split into orbits of sizes \{8\}, \{24\}, \{32\}, \{8, 16\}, \{8, 96\}} & True & 25.1 \\
4 & \texttt{Fourier modes of the auxiliary torus: phi\_\allowbreak{}n = Exp[I q\_\allowbreak{}n.y]/\allowbreak{}Sqrt[V\_\allowbreak{}perp] with q\_\allowbreak{}n = 2 Pi n/\allowbreak{}Lc (n in Z\textasciicircum{}4) are orthonormal for the proper volume element b\textasciicircum{}4 d\textasciicircum{}4y (exact integrals: n = m = (1, 0, -2, 3); n = (1, 0, -2, 3), m = (1, 1, -2, 3); n = 0, m = (0, 2, 0, 0))} & True & 25.2 \\
5 & \texttt{Parseval on the auxiliary torus: f(y) = Prod\_\allowbreak{}k cos\textasciicircum{}2(Pi y\_\allowbreak{}k/\allowbreak{}Lc) = Prod\_\allowbreak{}k (1/\allowbreak{}2 + (1/\allowbreak{}4) e\textasciicircum{}(2 Pi I y\_\allowbreak{}k/\allowbreak{}Lc) + (1/\allowbreak{}4) e\textasciicircum{}(-2 Pi I y\_\allowbreak{}k/\allowbreak{}Lc)) has the 81 non-zero coefficients c\_\allowbreak{}n = Sqrt[V\_\allowbreak{}perp] Prod\_\allowbreak{}k w(n\_\allowbreak{}k), w(0) = 1/\allowbreak{}2, w(+-1) = 1/\allowbreak{}4 (three of them checked by integration), and Sum\_\allowbreak{}n |c\_\allowbreak{}n|\textasciicircum{}2 == Int b\textasciicircum{}4 d\textasciicircum{}4y f\textasciicircum{}2 == (3/\allowbreak{}8)\textasciicircum{}4 V\_\allowbreak{}perp} & True & 25.2 \\
6 & \texttt{the Laplacian of the auxiliary torus: -(1/\allowbreak{}b\textasciicircum{}2) Sum\_\allowbreak{}k d\textasciicircum{}2 phi\_\allowbreak{}n/\allowbreak{}dy\_\allowbreak{}k\textasciicircum{}2 == (|q\_\allowbreak{}n|\textasciicircum{}2/\allowbreak{}b\textasciicircum{}2) phi\_\allowbreak{}n (n = (2, -1, 0, 3))} & True & 25.2 \\
7 & \texttt{scalar Kaluza-Klein masses: for the static metric diag(1, -1, -1, -1, -b0\textasciicircum{}2, ..., -b0\textasciicircum{}2) the Klein-Gordon operator Box + m\textasciicircum{}2 (Box f = |g|\textasciicircum{}(-1/\allowbreak{}2) d\_\allowbreak{}A(|g|\textasciicircum{}(1/\allowbreak{}2) g\textasciicircum{}AB d\_\allowbreak{}B f)) on Exp[-I E x0 + I k.x + I q.y] gives (-E\textasciicircum{}2 + |k|\textasciicircum{}2 + |q|\textasciicircum{}2/\allowbreak{}b0\textasciicircum{}2 + m\textasciicircum{}2) times the wave: in ordinary spacetime a particle of mass Sqrt[m\textasciicircum{}2 + |q|\textasciicircum{}2/\allowbreak{}b0\textasciicircum{}2]} & True & 25.2 \\
8 & \texttt{Jacobi's four-square theorem: the number of n in Z\textasciicircum{}4 with |n|\textasciicircum{}2 = k is r\_\allowbreak{}4(k) = 8 sigma(k) - 32 sigma(k/\allowbreak{}4) for k = 1..100 (direct count over \{-10, ..., 10\}\textasciicircum{}4)} & True & 25.3 \\
9 & \texttt{Weyl's law: the number N(R) of n in Z\textasciicircum{}4 with |n| <= R approaches the volume Pi\textasciicircum{}2 R\textasciicircum{}4/\allowbreak{}2 of the four-ball: N(10) = 49689, N(100) = 493490641, and |N(R)/\allowbreak{}(Pi\textasciicircum{}2 R\textasciicircum{}4/\allowbreak{}2) - 1| < 10\textasciicircum{}-2, 10\textasciicircum{}-3, 3 10\textasciicircum{}-5, 5 10\textasciicircum{}-6, 10\textasciicircum{}-6 for R = 10, 30, 100, 300, 1000} & True & 25.3 \\
10 & \texttt{the auxiliary torus today (M8 = 10 TeV, M\_\allowbreak{}Pl = 2.435 10\textasciicircum{}18 GeV): V\_\allowbreak{}perp = M\_\allowbreak{}Pl\textasciicircum{}2/\allowbreak{}M8\textasciicircum{}6 = 5.929 10\textasciicircum{}12 GeV\textasciicircum{}-4, l\_\allowbreak{}perp = V\_\allowbreak{}perp\textasciicircum{}(1/\allowbreak{}4) = b0 Lc = 1560.4 GeV\textasciicircum{}-1 = 3.079 10\textasciicircum{}-13 m, lowest Kaluza-Klein mass 2 Pi/\allowbreak{}l\_\allowbreak{}perp = 4.027 MeV; N(mu) = V\_\allowbreak{}perp mu\textasciicircum{}4/\allowbreak{}(32 Pi\textasciicircum{}2) = 1.17 10\textasciicircum{}13 levels below m\_\allowbreak{}r/\allowbreak{}2 = 5 GeV (the lattice count of section 22.2 agrees to 10\textasciicircum{}-5) and N(M8) = M\_\allowbreak{}Pl\textasciicircum{}2/\allowbreak{}(32 Pi\textasciicircum{}2 M8\textasciicircum{}2)} & True & 25.3 \\
11 & \texttt{Newton's law in eight dimensions: Omega\_\allowbreak{}6 = 2 Pi\textasciicircum{}(7/\allowbreak{}2)/\allowbreak{}Gamma(7/\allowbreak{}2) = 16 Pi\textasciicircum{}3/\allowbreak{}15, r\textasciicircum{}-5 is harmonic away from r = 0 in seven space dimensions ((1/\allowbreak{}r\textasciicircum{}6) d/\allowbreak{}dr (r\textasciicircum{}6 d/\allowbreak{}dr r\textasciicircum{}-5) = 0), and grad\textasciicircum{}2 Phi = 8 Pi G8 ((D - 3)/\allowbreak{}(D - 2)) M delta\textasciicircum{}7 (D = 8) gives Phi = -(5/\allowbreak{}6) kappa M/\allowbreak{}(5 Omega\_\allowbreak{}6 r\textasciicircum{}5) = -5 kappa M/\allowbreak{}(32 Pi\textasciicircum{}3 r\textasciicircum{}5); the same formula in D = 4 gives -G M/\allowbreak{}r} & True & 25.4 \\
12 & \texttt{images = Kaluza-Klein modes (Poisson summation; lengths in units of l\_\allowbreak{}perp): Sum\_\allowbreak{}n (r\textasciicircum{}2 + |n|\textasciicircum{}2)\textasciicircum{}(-5/\allowbreak{}2) == (4 Pi\textasciicircum{}2/\allowbreak{}(3 r)) Sum\_\allowbreak{}m Exp[-2 Pi r |m|] to 10\textasciicircum{}-8 at r = 1/\allowbreak{}4, 1/\allowbreak{}2, 1, 2 (left: shells with r\_\allowbreak{}4(k) to k = 1.5 10\textasciicircum{}6 and the continuum tail Pi\textasciicircum{}2 Int k (r\textasciicircum{}2 + k)\textasciicircum{}(-5/\allowbreak{}2) dk; right: shells to k = 20000)} & True & 25.4 \\
13 & \texttt{near and far field: Z(5/\allowbreak{}2) = Sum\_\allowbreak{}(n != 0) |n|\textasciicircum{}-5 = Sum\_\allowbreak{}k r\_\allowbreak{}4(k) k\textasciicircum{}(-5/\allowbreak{}2) = 7 zeta(5/\allowbreak{}2) zeta(3/\allowbreak{}2) = 24.531 (shell sum with tail, to 10\textasciicircum{}-8); near the source r\textasciicircum{}5 S(r) - 1 -> r\textasciicircum{}5 Z(5/\allowbreak{}2) (the ratio lies between 0.98 and 1 at r = 1/\allowbreak{}10 and 1/\allowbreak{}20, closer to 1 at 1/\allowbreak{}20): the eight-dimensional law; far from it 3 r S/\allowbreak{}(4 Pi\textasciicircum{}2) - 1 -> 8 Exp[-2 Pi r], the 8 lightest Kaluza-Klein modes (ratio 1 to 2 10\textasciicircum{}-2 at r = 2 and 3)} & True & 25.4 \\
14 & \texttt{graviton + radion: far from the source the images add up to (5 kappa M/\allowbreak{}(32 Pi\textasciicircum{}3)) (4 Pi\textasciicircum{}2/\allowbreak{}(3 V\_\allowbreak{}perp r)) == (5/\allowbreak{}3) G4 M/\allowbreak{}r with G4 = kappa/\allowbreak{}(8 Pi V\_\allowbreak{}perp); the radion of section 20.3, coupled to matter by (phi/\allowbreak{}(Sqrt[3] M\_\allowbreak{}Pl)) T, gives the potential -(1/\allowbreak{}(Sqrt[3] M\_\allowbreak{}Pl))\textasciicircum{}2 M M' Exp[-m\_\allowbreak{}r r]/\allowbreak{}(4 Pi r), and (1/\allowbreak{}(Sqrt[3] M\_\allowbreak{}Pl))\textasciicircum{}2/\allowbreak{}(4 Pi) == (2/\allowbreak{}3) G4 (8 Pi G4 = 1/\allowbreak{}M\_\allowbreak{}Pl\textasciicircum{}2): graviton 1 + massless radion 2/\allowbreak{}3 = 5/\allowbreak{}3} & True & 25.4 \\
15 & \texttt{Kaluza-Klein momentum is conserved in the bulk and not at the brane: Int b\textasciicircum{}4 d\textasciicircum{}4y phi\_\allowbreak{}n phi\_\allowbreak{}m phi\_\allowbreak{}l\textasciicircum{}* == 1/\allowbreak{}Sqrt[V\_\allowbreak{}perp] for l = n + m and 0 otherwise (exact: n = (1, 0, 0, 2), m = (0, 1, -1, 0), l = (1, 1, -1, 2) and (1, 1, -1, 1)); at y = 0, phi\_\allowbreak{}n phi\_\allowbreak{}m phi\_\allowbreak{}l\textasciicircum{}* == V\_\allowbreak{}perp\textasciicircum{}(-3/\allowbreak{}2) for every n, m, l} & True & 25.5 \\
16 & \texttt{the delta function of the torus: Int\_\allowbreak{}(-Lc/\allowbreak{}2)\textasciicircum{}(Lc/\allowbreak{}2) Exp[-I q\_\allowbreak{}n y] delta(y) dy == 1 for every integer n, so delta\textasciicircum{}4(y) = Lc\textasciicircum{}-4 Sum\_\allowbreak{}n Exp[I q\_\allowbreak{}n.y] (every mode with the weight 1/\allowbreak{}Lc\textasciicircum{}4); a brane source rho4 delta\textasciicircum{}4(y)/\allowbreak{}b\textasciicircum{}4 has Int b\textasciicircum{}4 d\textasciicircum{}4y (rho4 delta\textasciicircum{}4(y)/\allowbreak{}b\textasciicircum{}4) == rho4, i.e. the average rho4/\allowbreak{}V\_\allowbreak{}perp over the torus} & True & 25.5 \\
17 & \texttt{the world brane W = \{x4 = x5 = x6 = x7 = 0\}: the induced metric h = J\textasciicircum{}T.g.J of the embedding X = (x0, x1, x2, x3, 0, 0, 0, 0) is, for the cosmological metric of section 11.2, the Friedmann-Robertson-Walker metric diag(1, -a\textasciicircum{}2, -a\textasciicircum{}2, -a\textasciicircum{}2) (independent of b); for the inhomogeneous metric of section 11.3 it is diag(n\textasciicircum{}2, -a\textasciicircum{}2, -a\textasciicircum{}2, -a\textasciicircum{}2) at x4 = 0 with n = 1 + x0\textasciicircum{}2/\allowbreak{}11, a = 1 + x0/\allowbreak{}3} & True & 25.6 \\
18 & \texttt{brane integrals in eight-dimensional form: at the brane Sqrt|det g| == Sqrt|det h| b\textasciicircum{}4 with b\textasciicircum{}4 = Sqrt[det(-g\_\allowbreak{}kl)] (cosmological metric; inhomogeneous metric exactly at x0 = 1/\allowbreak{}2), so Int\_\allowbreak{}W d\textasciicircum{}4x Sqrt|h| L == Int d\textasciicircum{}8X Sqrt|g| (delta\textasciicircum{}4(y)/\allowbreak{}b\textasciicircum{}4) L; test: Int d\textasciicircum{}4y b\textasciicircum{}4 (delta\textasciicircum{}4(y)/\allowbreak{}b\textasciicircum{}4) F(y) == F(0) = 1 for F = Exp[x4 + 2 x5] Cos[x6 - x7]} & True & 25.6 \\
19 & \texttt{how the brane sits in the bulk: the Christoffel symbols \{k mu nu\} at the brane (k = 4..7, mu, nu = 0..3; = (1/\allowbreak{}(2 b\textasciicircum{}2)) d\_\allowbreak{}k g\_\allowbreak{}mu nu for these metrics, b times the extrinsic curvature) all vanish for the cosmological metric: the brane is totally geodesic, and a free particle moving along it stays on it. For the inhomogeneous metric at x0 = 1/\allowbreak{}2 the only non-zero one is \{4 00\} = n (d\_\allowbreak{}4 n)/\allowbreak{}b\textasciicircum{}2 = 1458/\allowbreak{}58619: a free particle at rest on the brane is accelerated off it} & True & 25.6 \\
20 & \texttt{explicit forms: every Gamma17[A] is a signed permutation matrix (one entry +-1 in each row and column), and 2 Sqrt[2] U has only the entries 0, +-1, +-I, eight non-zero ones in each row and in each column} & True & 25.7 \\
21 & \texttt{brane spinor (x) transverse spinor: Gamma17[mu].Gamma17[nu] == U.(gamma\_\allowbreak{}mu.gamma\_\allowbreak{}nu (x) 1\_\allowbreak{}4).U\textasciicircum{}dag, Gamma17[k].Gamma17[l] == U.(1\_\allowbreak{}4 (x) rho\_\allowbreak{}k.rho\_\allowbreak{}l).U\textasciicircum{}dag and Gamma17[mu].Gamma17[k] == U.(gamma\_\allowbreak{}mu.gamma5 (x) rho\_\allowbreak{}k).U\textasciicircum{}dag (mu, nu = 0..3, k, l = 4..7): the Lorentz group of the brane acts on the brane spinor index, the rotations about the brane on the transverse index, and the generators that mix them are broken by the brane} & True & 25.7 \\
22 & \texttt{the transverse index carries Spin(4) = SU(2) x SU(2): rho4567 == -gamma5 (4 x 4), it commutes with every rho\_\allowbreak{}k.rho\_\allowbreak{}l and anticommutes with every rho\_\allowbreak{}k; for each even permutation (k, l, m, n) of (4, 5, 6, 7) the combination rho\_\allowbreak{}k.rho\_\allowbreak{}l + s rho\_\allowbreak{}m.rho\_\allowbreak{}n vanishes on the eigenspace rho4567 = s, and rho\_\allowbreak{}k.rho\_\allowbreak{}l - s rho\_\allowbreak{}m.rho\_\allowbreak{}n acts there with rank 2: the four values of a form (2, 1) + (1, 2)} & True & 25.7 \\
23 & \texttt{chirality at the brane: PL17 == U.(P\_\allowbreak{}L (x) P\_\allowbreak{}+ + P\_\allowbreak{}R (x) P\_\allowbreak{}-).U\textasciicircum{}dag with P\_\allowbreak{}L,R = (1 -+ gamma5)/\allowbreak{}2 and P\_\allowbreak{}+- = (1 +- rho4567)/\allowbreak{}2: the left-handed half psi\_\allowbreak{}L (chi17 = -1, 8 components) is two brane-left-handed Weyl spinors with rho4567 = +1 and two brane-right-handed ones with rho4567 = -1 (each part of rank 4)} & True & 25.7 \\
24 & \texttt{an explicit spinor: psi = U.(u (x) e\_\allowbreak{}1) with the brane spinor u = (1, 2, 0, -1): psibar psi == ubar u == 4 and psibar Gamma\textasciicircum{}mu psi == ubar gamma\textasciicircum{}mu u (mu = 0..3)} & True & 25.7 \\
25 & \texttt{brane tension: for S\_\allowbreak{}W = -Int d\textasciicircum{}4x Sqrt|h| f\_\allowbreak{}W and h = diag(1/\allowbreak{}u0, -1/\allowbreak{}u1, -1/\allowbreak{}u2, -1/\allowbreak{}u3) (h\textasciicircum{}00 = u0, h\textasciicircum{}ii = -u\_\allowbreak{}i), T\_\allowbreak{}mu nu = (2/\allowbreak{}Sqrt|h|) dS\_\allowbreak{}W/\allowbreak{}dh\textasciicircum{}mu nu == f\_\allowbreak{}W h\_\allowbreak{}mu nu, T\textasciicircum{}0\_\allowbreak{}0 == f\_\allowbreak{}W (energy density) and T\textasciicircum{}1\_\allowbreak{}1 == f\_\allowbreak{}W (pressure -f\_\allowbreak{}W: w = -1 along the brane)} & True & 25.8 \\
26 & \texttt{sources smeared over the torus (W(a, b) the comoving energy, section 20.2): the brane tension W = a\textasciicircum{}3 f\_\allowbreak{}W/\allowbreak{}Lc\textasciicircum{}4 gives eps = f\_\allowbreak{}W/\allowbreak{}V\_\allowbreak{}perp, pa = -eps, pb = 0; brane radiation W = a\textasciicircum{}3 r0 a\textasciicircum{}-4/\allowbreak{}Lc\textasciicircum{}4 gives pa = eps/\allowbreak{}3, pb = 0; radiation of the bulk, isotropic in the seven space directions, has <cos\textasciicircum{}2 theta> = 3/\allowbreak{}7 and <sin\textasciicircum{}2 theta> = 4/\allowbreak{}7 for the weight cos\textasciicircum{}2 theta sin\textasciicircum{}3 theta of the momentum sphere, so pa = (1/\allowbreak{}3)(3/\allowbreak{}7) eps = eps/\allowbreak{}7 and pb = (1/\allowbreak{}4)(4/\allowbreak{}7) eps = eps/\allowbreak{}7} & True & 25.8 \\
27 & \texttt{the brane in four-dimensional physics: (i) the tension is vacuum energy, kappa f\_\allowbreak{}W/\allowbreak{}V\_\allowbreak{}perp == f\_\allowbreak{}W/\allowbreak{}M\_\allowbreak{}Pl\textasciicircum{}2 (V\_\allowbreak{}perp = kappa M\_\allowbreak{}Pl\textasciicircum{}2), and f\_\allowbreak{}W = M8\textasciicircum{}4 is between 3 10\textasciicircum{}62 and 5 10\textasciicircum{}62 times rho\_\allowbreak{}Lambda = 3 H0\textasciicircum{}2 M\_\allowbreak{}Pl\textasciicircum{}2 Omega\_\allowbreak{}Lambda (section 24.5); (ii) normalcy temperature: (T\textasciicircum{}7/\allowbreak{}M8\textasciicircum{}6)/\allowbreak{}(T\textasciicircum{}2/\allowbreak{}M\_\allowbreak{}Pl) == 1 at T* = (M8\textasciicircum{}6/\allowbreak{}M\_\allowbreak{}Pl)\textasciicircum{}(1/\allowbreak{}5) = 13.3 GeV, and V\_\allowbreak{}perp T\textasciicircum{}4 T\textasciicircum{}3/\allowbreak{}M\_\allowbreak{}Pl\textasciicircum{}2 == T\textasciicircum{}7/\allowbreak{}M8\textasciicircum{}6; (iii) thin brane: (1/\allowbreak{}y\textasciicircum{}3) d/\allowbreak{}dy (y\textasciicircum{}3 d/\allowbreak{}dy y\textasciicircum{}-2) == 0 and 2 Pi\textasciicircum{}2 y\textasciicircum{}3 d/\allowbreak{}dy y\textasciicircum{}-2 == -4 Pi\textasciicircum{}2 (the Green's function -1/\allowbreak{}(4 Pi\textasciicircum{}2 |y|\textasciicircum{}2) of the four transverse directions), and Sqrt[kappa f\_\allowbreak{}W] == 1/\allowbreak{}M8 for f\_\allowbreak{}W = M8\textasciicircum{}4} & True & 25.8 \\*
\end{longtable}}
The Python program’s boolean checks:

{\footnotesize
\begin{longtable}{>{\raggedright\arraybackslash}p{\dimexpr 0.0500\linewidth-2\tabcolsep\relax}>{\raggedright\arraybackslash}p{\dimexpr 0.8300\linewidth-2\tabcolsep\relax}>{\raggedright\arraybackslash}p{\dimexpr 0.1200\linewidth-2\tabcolsep\relax}}
\toprule
\textbf{\#} & \textbf{check} & \textbf{result} \\
\midrule\endfirsthead
\toprule
\textbf{\#} & \textbf{check} & \textbf{result} \\
\midrule\endhead
\endfoot
\bottomrule\endlastfoot
1 & \texttt{torus: det(-g\_\allowbreak{}kl) = b\textasciicircum{}8 for the cosmological metric; the proper volume of [0, Lc)\textasciicircum{}4 is b\textasciicircum{}4 Lc\textasciicircum{}4} & True \\
2 & \texttt{only b Lc is physical: y = lambda y' maps (Lc, b) to (Lc/\allowbreak{}lambda, lambda b); b\textasciicircum{}2 dy.dy, b Lc and b\textasciicircum{}4 Lc\textasciicircum{}4 unchanged} & True \\
3 & \texttt{symmetries fixing y = 0: 384 = 2\textasciicircum{}4 4! distinct signed permutation matrices, orthogonal integer matrices; shells 1, 2, 3, 4, 9 have orbits \{8\}, \{24\}, \{32\}, \{8, 16\}, \{8, 96\}} & True \\
4 & \texttt{Fourier modes phi\_\allowbreak{}n = exp(i q\_\allowbreak{}n.y)/\allowbreak{}sqrt(V\_\allowbreak{}perp), q\_\allowbreak{}n = 2 pi n/\allowbreak{}Lc: orthonormal for b\textasciicircum{}4 d\textasciicircum{}4y (three exact pairs)} & True \\
5 & \texttt{Parseval: f = prod cos\textasciicircum{}2(pi y\_\allowbreak{}k/\allowbreak{}Lc) has 81 coefficients sqrt(V\_\allowbreak{}perp) prod w(n\_\allowbreak{}k) (3 checked by integration), Sum |c\_\allowbreak{}n|\textasciicircum{}2 = Int b\textasciicircum{}4 f\textasciicircum{}2 = (3/\allowbreak{}8)\textasciicircum{}4 V\_\allowbreak{}perp} & True \\
6 & \texttt{Laplacian of the torus: -(1/\allowbreak{}b\textasciicircum{}2) Sum d\textasciicircum{}2/\allowbreak{}dy\_\allowbreak{}k\textasciicircum{}2 phi\_\allowbreak{}n = (|q\_\allowbreak{}n|\textasciicircum{}2/\allowbreak{}b\textasciicircum{}2) phi\_\allowbreak{}n} & True \\
7 & \texttt{scalar Kaluza-Klein masses: (Box + m\textasciicircum{}2) exp(-iEx0 + ik.x + iq.y) = (-E\textasciicircum{}2 + |k|\textasciicircum{}2 + |q|\textasciicircum{}2/\allowbreak{}b0\textasciicircum{}2 + m\textasciicircum{}2) exp(...) on diag(1,-1,-1,-1,-b0\textasciicircum{}2,...)} & True \\
8 & \texttt{graviton on the torus: D(D-3)/\allowbreak{}2 = 20 (D = 8) = 2 + 8 + 10 (zero mode: graviton, 4 graviphotons, 10 moduli) = 5 + 9 + 6 (massive level)} & True \\
9 & \texttt{Jacobi: \#\{n in Z\textasciicircum{}4 : |n|\textasciicircum{}2 = k\} = r\_\allowbreak{}4(k) = 8 sigma(k) - 32 sigma(k/\allowbreak{}4), direct count for k = 1..100} & True \\
10 & \texttt{Weyl's law: N(10) = 49689, N(100) = 493490641; |N(R)/\allowbreak{}(pi\textasciicircum{}2 R\textasciicircum{}4/\allowbreak{}2) - 1| < 1e-2, 1e-3, 3e-5, 5e-6, 1e-6} & True \\
11 & \texttt{numbers today: V\_\allowbreak{}perp = 5.929e12 GeV\textasciicircum{}-4, l\_\allowbreak{}perp = 1560.4 GeV\textasciicircum{}-1 = 3.079e-13 m, 2 pi/\allowbreak{}l\_\allowbreak{}perp = 4.027 MeV; N(5 GeV) = 1.17e13 (lattice to 1e-5); N(M8) = M\_\allowbreak{}Pl\textasciicircum{}2/\allowbreak{}(32 pi\textasciicircum{}2 M8\textasciicircum{}2)} & True \\
12 & \texttt{from eight to four dimensions: for diag(1, -a\textasciicircum{}2, -a\textasciicircum{}2, -a\textasciicircum{}2, -b0\textasciicircum{}2, ..., -b0\textasciicircum{}2) (static torus) R = R4 = -6(Ha' + 2 Ha\textasciicircum{}2), the scalar curvature of the FRW metric diag(1, -a\textasciicircum{}2, -a\textasciicircum{}2, -a\textasciicircum{}2); sqrt|g| = a\textasciicircum{}3 b0\textasciicircum{}4} & True \\
13 & \texttt{quaternion subalgebra of the brane directions: e1 e2 = -e3, e2 e3 = -e1, e3 e1 = -e2 (e\_\allowbreak{}i e\_\allowbreak{}j = -delta\_\allowbreak{}ij - eps\_\allowbreak{}ijk e\_\allowbreak{}k), span\{e0, ..., e3\} closed under the octonion product, Re(e\_\allowbreak{}A e\_\allowbreak{}B) = diag(1, -1, -1, -1)} & True \\
14 & \texttt{Newton in 8D: Omega\_\allowbreak{}6 = 16 pi\textasciicircum{}3/\allowbreak{}15, (1/\allowbreak{}r\textasciicircum{}6) d/\allowbreak{}dr (r\textasciicircum{}6 d/\allowbreak{}dr r\textasciicircum{}-5) = 0, Phi = -(5/\allowbreak{}6) kappa M/\allowbreak{}(5 Omega\_\allowbreak{}6 r\textasciicircum{}5) = -5 kappa M/\allowbreak{}(32 pi\textasciicircum{}3 r\textasciicircum{}5); D = 4: -G M/\allowbreak{}r} & True \\
15 & \texttt{4D Fourier transform: Int d\textasciicircum{}4y exp(-ip.y) (r\textasciicircum{}2 + |y|\textasciicircum{}2)\textasciicircum{}(-5/\allowbreak{}2) = (4 pi\textasciicircum{}2/\allowbreak{}(3 r)) exp(-p r) (mpmath, three (p, r), to 1e-20)} & True \\
16 & \texttt{Poisson duality: Sum\_\allowbreak{}n (r\textasciicircum{}2 + |n|\textasciicircum{}2)\textasciicircum{}(-5/\allowbreak{}2) = (4 pi\textasciicircum{}2/\allowbreak{}(3 r)) Sum\_\allowbreak{}m exp(-2 pi r |m|) to 1e-8 at r = 1/\allowbreak{}4, 1/\allowbreak{}2, 1, 2} & True \\
17 & \texttt{Z(5/\allowbreak{}2) = 7 zeta(5/\allowbreak{}2) zeta(3/\allowbreak{}2) = Sum r\_\allowbreak{}4(k) k\textasciicircum{}-5/\allowbreak{}2 (to 1e-8); near field ratio -> 1 (0.98 < . < 1 at r = 1/\allowbreak{}10, 1/\allowbreak{}20); far field ratio -> 1 (to 2e-2 at r = 2, 3)} & True \\
18 & \texttt{graviton + radion: (5 kappa M/\allowbreak{}(32 pi\textasciicircum{}3))(4 pi\textasciicircum{}2/\allowbreak{}(3 V\_\allowbreak{}perp r)) = (5/\allowbreak{}3) G4 M/\allowbreak{}r; radion (1/\allowbreak{}(sqrt3 M\_\allowbreak{}Pl))\textasciicircum{}2/\allowbreak{}(4 pi) = (2/\allowbreak{}3)/\allowbreak{}(8 pi M\_\allowbreak{}Pl\textasciicircum{}2)} & True \\
19 & \texttt{Kaluza-Klein momentum: bulk vertex Int b\textasciicircum{}4 phi\_\allowbreak{}n phi\_\allowbreak{}m phi\_\allowbreak{}l\textasciicircum{}* = 1/\allowbreak{}sqrt(V\_\allowbreak{}perp) for l = n + m, 0 otherwise; at y = 0 phi\_\allowbreak{}n phi\_\allowbreak{}m phi\_\allowbreak{}l\textasciicircum{}* = V\_\allowbreak{}perp\textasciicircum{}(-3/\allowbreak{}2)} & True \\
20 & \texttt{delta function of the torus: Fourier coefficients Int exp(-i q\_\allowbreak{}n y) delta(y) dy = 1; Int b\textasciicircum{}4 d\textasciicircum{}4y rho4 delta\textasciicircum{}4(y)/\allowbreak{}b\textasciicircum{}4 = rho4 (average rho4/\allowbreak{}V\_\allowbreak{}perp)} & True \\
21 & \texttt{induced metric h = J\textasciicircum{}T g J: FRW diag(1, -a\textasciicircum{}2, -a\textasciicircum{}2, -a\textasciicircum{}2) for the cosmological metric; diag(n\textasciicircum{}2, -a\textasciicircum{}2, -a\textasciicircum{}2, -a\textasciicircum{}2) at x4 = 0 with n = 1 + x0\textasciicircum{}2/\allowbreak{}11, a = 1 + x0/\allowbreak{}3 for the inhomogeneous metric} & True \\
22 & \texttt{brane volume element: sqrt|g| = sqrt|h| b\textasciicircum{}4 at the brane (cosmological metric; inhomogeneous metric at x0 = 1/\allowbreak{}2, exact), so Int\_\allowbreak{}W d\textasciicircum{}4x sqrt|h| L = Int d\textasciicircum{}8X sqrt|g| delta\textasciicircum{}4(y)/\allowbreak{}b\textasciicircum{}4 L; test Int b\textasciicircum{}4 (delta\textasciicircum{}4/\allowbreak{}b\textasciicircum{}4) F = F(0)} & True \\
23 & \texttt{Christoffel symbols \{k mu nu\} at the brane: all 0 for the cosmological metric (totally geodesic); for the inhomogeneous metric at x0 = 1/\allowbreak{}2 only \{4 00\} = 1458/\allowbreak{}58619 = n d4n/\allowbreak{}b\textasciicircum{}2} & True \\
24 & \texttt{R\_\allowbreak{}A = (gamma\_\allowbreak{}mu (x) 1\_\allowbreak{}4, gamma5 (x) rho\_\allowbreak{}k) and Gamma17[A] both satisfy X\_\allowbreak{}A X\_\allowbreak{}B + X\_\allowbreak{}B X\_\allowbreak{}A = 2 eta\_\allowbreak{}AB 1\_\allowbreak{}16} & True \\
25 & \texttt{every Gamma17[A] is a signed permutation matrix} & True \\
26 & \texttt{intertwiner: Schur average with the seed matrix unit (1, 2); 2 sqrt(2) U has entries 0, +-1, +-i (rounding error < 1e-12) with 8 non-zero entries in each row and column; Gamma17[A].U8 = U8.R\_\allowbreak{}A and U8\textasciicircum{}dag.U8 = 8 1\_\allowbreak{}16 exactly} & True \\
27 & \texttt{Schur's lemma: the linear map X -> (Gamma\_\allowbreak{}A X - X R\_\allowbreak{}A)\_\allowbreak{}A has a one-dimensional null space (one singular value < 1e-10, the next > 0.1)} & True \\
28 & \texttt{brane Lorentz group on the brane spinor index, transverse rotations on a: Gamma\_\allowbreak{}mu Gamma\_\allowbreak{}nu = U (gamma\_\allowbreak{}mu gamma\_\allowbreak{}nu (x) 1) U\textasciicircum{}dag, Gamma\_\allowbreak{}k Gamma\_\allowbreak{}l = U (1 (x) rho\_\allowbreak{}k rho\_\allowbreak{}l) U\textasciicircum{}dag, Gamma\_\allowbreak{}mu Gamma\_\allowbreak{}k = U (gamma\_\allowbreak{}mu gamma5 (x) rho\_\allowbreak{}k) U\textasciicircum{}dag} & True \\
29 & \texttt{Spin(4) on the transverse index: rho4567 = -gamma5 (4x4) commutes with rho\_\allowbreak{}k rho\_\allowbreak{}l, anticommutes with rho\_\allowbreak{}k; (rho\_\allowbreak{}k rho\_\allowbreak{}l + s rho\_\allowbreak{}m rho\_\allowbreak{}n) P\_\allowbreak{}s = 0 and (rho\_\allowbreak{}k rho\_\allowbreak{}l - s rho\_\allowbreak{}m rho\_\allowbreak{}n) P\_\allowbreak{}s has rank 2 for even (k,l,m,n): 4 = (2,1) + (1,2)} & True \\
30 & \texttt{chirality: chi = i Gamma\_\allowbreak{}0...Gamma\_\allowbreak{}7 = U (gamma5 (x) rho4567) U\textasciicircum{}dag; P\_\allowbreak{}L17 = U (P\_\allowbreak{}L (x) P\_\allowbreak{}+ + P\_\allowbreak{}R (x) P\_\allowbreak{}-) U\textasciicircum{}dag, each part of rank 4} & True \\
31 & \texttt{example psi = U.(u (x) e\_\allowbreak{}1), u = (1, 2, 0, -1): psibar psi = ubar u = 4 and psibar Gamma\textasciicircum{}mu psi = ubar gamma\textasciicircum{}mu u} & True \\
32 & \texttt{four brane Dirac spinors: psibar psi' = Sum\_\allowbreak{}a psibar\_\allowbreak{}a psi'\_\allowbreak{}a and psibar Gamma\textasciicircum{}mu psi' = Sum\_\allowbreak{}a psibar\_\allowbreak{}a gamma\textasciicircum{}mu psi'\_\allowbreak{}a (random integer spinors)} & True \\
33 & \texttt{Kaluza-Klein operator X = gamma5 (x) (rho.q)/\allowbreak{}b0: X\textasciicircum{}2 = -|q|\textasciicircum{}2/\allowbreak{}b0\textasciicircum{}2 1\_\allowbreak{}16 and X anticommutes with gamma\_\allowbreak{}mu (x) 1\_\allowbreak{}4} & True \\
34 & \texttt{brane tension: T\_\allowbreak{}mu nu = (2/\allowbreak{}sqrt|h|) dS/\allowbreak{}dh\textasciicircum{}mu nu = f\_\allowbreak{}W h\_\allowbreak{}mu nu for S = -Int sqrt|h| f\_\allowbreak{}W; T\textasciicircum{}0\_\allowbreak{}0 = f\_\allowbreak{}W, T\textasciicircum{}1\_\allowbreak{}1 = f\_\allowbreak{}W (p = -f\_\allowbreak{}W, w = -1)} & True \\
35 & \texttt{smeared sources: tension eps = f\_\allowbreak{}W/\allowbreak{}V\_\allowbreak{}perp, pa = -eps, pb = 0; brane radiation pa = eps/\allowbreak{}3, pb = 0; bulk radiation <cos\textasciicircum{}2> = 3/\allowbreak{}7, <sin\textasciicircum{}2> = 4/\allowbreak{}7: pa = pb = eps/\allowbreak{}7} & True \\
36 & \texttt{normalcy temperature: (T\textasciicircum{}7/\allowbreak{}M8\textasciicircum{}6)/\allowbreak{}(T\textasciicircum{}2/\allowbreak{}M\_\allowbreak{}Pl) = 1 at T* = (M8\textasciicircum{}6/\allowbreak{}M\_\allowbreak{}Pl)\textasciicircum{}(1/\allowbreak{}5) = 13.3 GeV; V\_\allowbreak{}perp T\textasciicircum{}4 T\textasciicircum{}3/\allowbreak{}M\_\allowbreak{}Pl\textasciicircum{}2 = T\textasciicircum{}7/\allowbreak{}M8\textasciicircum{}6} & True \\
37 & \texttt{the tension is vacuum energy: kappa f\_\allowbreak{}W/\allowbreak{}V\_\allowbreak{}perp = f\_\allowbreak{}W/\allowbreak{}M\_\allowbreak{}Pl\textasciicircum{}2 (V\_\allowbreak{}perp = kappa M\_\allowbreak{}Pl\textasciicircum{}2); f\_\allowbreak{}W = M8\textasciicircum{}4 is 3e62 to 5e62 rho\_\allowbreak{}Lambda} & True \\
38 & \texttt{Green's function of the four transverse directions: (1/\allowbreak{}y\textasciicircum{}3) d/\allowbreak{}dy (y\textasciicircum{}3 d/\allowbreak{}dy y\textasciicircum{}-2) = 0, flux 2 pi\textasciicircum{}2 y\textasciicircum{}3 d/\allowbreak{}dy y\textasciicircum{}-2 = -4 pi\textasciicircum{}2; sqrt(kappa f\_\allowbreak{}W) = 1/\allowbreak{}M8 for f\_\allowbreak{}W = M8\textasciicircum{}4} & True \\*
\end{longtable}}
The numbers printed by both programs (largest relative difference per group; integers must agree exactly):

{\footnotesize
\begin{longtable}{>{\raggedright\arraybackslash}p{\dimexpr 0.5200\linewidth-2\tabcolsep\relax}>{\raggedright\arraybackslash}p{\dimexpr 0.1000\linewidth-2\tabcolsep\relax}>{\raggedright\arraybackslash}p{\dimexpr 0.2200\linewidth-2\tabcolsep\relax}>{\raggedright\arraybackslash}p{\dimexpr 0.1600\linewidth-2\tabcolsep\relax}}
\toprule
\textbf{quantity} & \textbf{values} & \textbf{max. relative difference} & \textbf{tolerance} \\
\midrule\endfirsthead
\toprule
\textbf{quantity} & \textbf{values} & \textbf{max. relative difference} & \textbf{tolerance} \\
\midrule\endhead
\endfoot
\bottomrule\endlastfoot
orbit sizes of the shells $k=1,\dots,12$ & 16 & 0 & exact \\
$r_4(k)$, $k=0,\dots,16$ & 17 & 0 & exact \\
Weyl: $N(R)$, $R=10,30,100,300,1000$ & 5 & 0 & exact \\
Weyl: $N(R)/(\pi^2R^4/2)-1$ & 5 & $3.8\times10^{-10}$ & $1\times10^{-9}$ \\
$V_\perp$, $\ell_\perp$ (GeV$^{-1}$, m, fm), $2\pi/\ell_\perp$ & 5 & $2.2\times10^{-16}$ & $1\times10^{-12}$ \\
$\ell_\perp/r_p$, $a_0/\ell_\perp$, $R=\ell_\perp/2\pi$, $M_D$ & 5 & $1.7\times10^{-16}$ & $1\times10^{-12}$ \\
$N(\mu)$ at 5 GeV, $m_h$, $M_8$; lattice count below 5 GeV & 4 & $1.8\times10^{-16}$ & $1\times10^{-12}$ \\
the first levels ($k$, $r_4$, $4r_4$, mass) & 32 & $2.5\times10^{-16}$ & $1\times10^{-12}$ \\
image sums $S(r)$ and Kaluza–Klein sums, $r=\frac14,\frac12,1,2$ & 8 & $1.2\times10^{-15}$ & $1\times10^{-9}$ \\
$Z(5/2)$: closed form, shell sum & 2 & $2.9\times10^{-16}$ & $1\times10^{-9}$ \\
near-field and far-field ratios & 5 & $7.8\times10^{-12}$ & $1\times10^{-6}$ \\
Christoffel symbol $\{4\,00\}$ at the brane & 1 & 0 & $1\times10^{-15}$ \\
$\Gamma_A$ as signed permutations & 128 & 0 & exact \\
$\mathrm{Re}(2\sqrt2\,U)$ & 256 & 0 & exact \\
$\mathrm{Im}(2\sqrt2\,U)$ & 256 & 0 & exact \\
example spinor: $\bar\psi\psi$, $\bar\psi\Gamma^\mu\psi$ & 5 & 0 & $1\times10^{-15}$ \\
$T_*$, $f_W/\rho_\Lambda$, $1/(2\pi M_8)$ & 3 & $5.8\times10^{-16}$ & $1\times10^{-12}$ \\*
\end{longtable}}
The other results of this document (Examples 7.1–7.3 and 10.1–10.2, sections 8 and 11) are outputs of sections 17–24 of 7+1.nb, evaluated in the same run; they rest on these checks of the notebook:

{\footnotesize
\begin{longtable}{>{\raggedright\arraybackslash}p{\dimexpr 0.0500\linewidth-2\tabcolsep\relax}>{\raggedright\arraybackslash}p{\dimexpr 0.7700\linewidth-2\tabcolsep\relax}>{\raggedright\arraybackslash}p{\dimexpr 0.0800\linewidth-2\tabcolsep\relax}>{\raggedright\arraybackslash}p{\dimexpr 0.1000\linewidth-2\tabcolsep\relax}}
\toprule
\textbf{\#} & \textbf{check} & \textbf{result} & \textbf{section} \\
\midrule\endfirsthead
\toprule
\textbf{\#} & \textbf{check} & \textbf{result} & \textbf{section} \\
\midrule\endhead
\endfoot
\bottomrule\endlastfoot
1 & \texttt{brane frame metric diag(+1, -1, -1, -1) == eta17 restricted to x0..x3: the squares of the octonion basis e[0], e[1], e[2], e[3] of the quaternion subalgebra (Re(e[A] e[B]), A, B = 0..3), which is closed under the octonion product} & True & 17 \\
2 & \texttt{anomaly cancellation per generation (with nu\_\allowbreak{}R): SU(3)\textasciicircum{}2U(1), SU(2)\textasciicircum{}2U(1), U(1)\textasciicircum{}3, grav\textasciicircum{}2U(1), SU(3)\textasciicircum{}3 all vanish; doublets per generation even} & True & 17.1 \\
3 & \texttt{electric charges Q = T3 + Y: (u, d) = (2/\allowbreak{}3, -1/\allowbreak{}3), (nu, e) = (0, -1), u\_\allowbreak{}R 2/\allowbreak{}3, d\_\allowbreak{}R -1/\allowbreak{}3, e\_\allowbreak{}R -1, nu\_\allowbreak{}R 0} & True & 17.1 \\
4 & \texttt{g* = 28 + (7/\allowbreak{}8) 90 = 106.75 for the Standard Model fields (without nu\_\allowbreak{}R)} & True & 17.1 \\
5 & \texttt{gauge-boson mass matrix: eigenvalues g\textasciicircum{}2 v\textasciicircum{}2/\allowbreak{}4 (twice), (g\textasciicircum{}2 + g'\textasciicircum{}2) v\textasciicircum{}2/\allowbreak{}4 and 0 (photon)} & True & 17.3 \\
6 & \texttt{tree level: m\_\allowbreak{}W == g v/\allowbreak{}2, m\_\allowbreak{}Z == Sqrt[g\textasciicircum{}2 + g'\textasciicircum{}2] v/\allowbreak{}2, rho == m\_\allowbreak{}W\textasciicircum{}2/\allowbreak{}(m\_\allowbreak{}Z\textasciicircum{}2 cos\textasciicircum{}2 theta\_\allowbreak{}W) == 1, m\_\allowbreak{}H\textasciicircum{}2 == 2 lambda v\textasciicircum{}2} & True & 17.3 \\
7 & \texttt{age of the universe from the brane Friedmann equation: 13.79 Gyr (Planck 2018: 13.787 +- 0.020 Gyr)} & True & 17.5 \\
8 & \texttt{ADD relation: kappa8 rho4/\allowbreak{}V\_\allowbreak{}perp == 8 Pi G4 rho4 with G4 = G8/\allowbreak{}V\_\allowbreak{}perp (kappa8 = 8 Pi G8)} & True & 17.5 \\
9 & \texttt{brane representation: R\_\allowbreak{}A R\_\allowbreak{}B + R\_\allowbreak{}B R\_\allowbreak{}A == 2 eta17\_\allowbreak{}AB ID16 for R = (gamma\_\allowbreak{}mu (x) 1\_\allowbreak{}4, gamma5 (x) rho\_\allowbreak{}k)} & True & 18 \\
10 & \texttt{intertwiner: Gamma17[A].U == U.R\_\allowbreak{}A for A = 0..7, and U is unitary} & True & 18 \\
11 & \texttt{chi17 == U.(gamma5 (x) rho4567).U\textasciicircum{}dag; rho4567 Hermitian with eigenvalues \{-1, -1, 1, 1\}} & True & 18 \\
12 & \texttt{bilinears at the brane: psibar psi' == Sum\_\allowbreak{}a psibar\_\allowbreak{}a psi'\_\allowbreak{}a and psibar Gamma\textasciicircum{}mu psi' == Sum\_\allowbreak{}a psibar\_\allowbreak{}a gamma\textasciicircum{}mu psi'\_\allowbreak{}a (mu = 0..3)} & True & 18 \\
13 & \texttt{KK operator: X.X == -|q|\textasciicircum{}2/\allowbreak{}b0\textasciicircum{}2 ID16 and X anticommutes with gamma\_\allowbreak{}mu (x) 1\_\allowbreak{}4} & True & 18.1 \\
14 & \texttt{KK masses: Det[gamma.p (x) 1 + X - m] vanishes on p0\textasciicircum{}2 - |p|\textasciicircum{}2 == m\textasciicircum{}2 + |q|\textasciicircum{}2/\allowbreak{}b0\textasciicircum{}2 (40-digit arithmetic) and not off shell} & True & 18.1 \\
15 & \texttt{portal mixing: 1 - N.N\textasciicircum{}dag == mu.mu\textasciicircum{}dag/\allowbreak{}m\textasciicircum{}2 to second order (non-unitary light mixing)} & True & 18.2 \\
16 & \texttt{brane tension smeared over the torus (V\_\allowbreak{}perp = b\textasciicircum{}4 Lc\textasciicircum{}4): eps = f\_\allowbreak{}W/\allowbreak{}V\_\allowbreak{}perp, pa = -f\_\allowbreak{}W/\allowbreak{}V\_\allowbreak{}perp, pb = 0} & True & 20.2 \\
17 & \texttt{brane fluid rho\_\allowbreak{}4 = N a\textasciicircum{}(-3(1+w)), smeared: eps = rho\_\allowbreak{}4/\allowbreak{}V\_\allowbreak{}perp, pa = w eps, pb = 0} & True & 20.2 \\
18 & \texttt{every such source is conserved: eps' + 3Ha(eps + pa) + 4Hb(eps + pb) == 0 for arbitrary W(a, b), a(x0), b(x0)} & True & 20.2 \\
19 & \texttt{brane Friedmann equation: 3 Ha\textasciicircum{}2 == (kappa/\allowbreak{}V\_\allowbreak{}perp) rho\_\allowbreak{}4 + Lambda4 with Lambda4 = Lambda + kappa U0 + kappa f\_\allowbreak{}W/\allowbreak{}V\_\allowbreak{}perp, i.e. 8 Pi G4 = kappa/\allowbreak{}V\_\allowbreak{}perp} & True & 20.3 \\
20 & \texttt{Einstein-frame form of the constraint: 3 (Ha + 2 Hb)\textasciicircum{}2 == 6 Hb\textasciicircum{}2 + (E\textasciicircum{}(1)0\_\allowbreak{}0 of section 12.6), i.e. 3 M\_\allowbreak{}Pl\textasciicircum{}2 H\_\allowbreak{}E\textasciicircum{}2 == phi'\textasciicircum{}2/\allowbreak{}2 + rho with phi = Sqrt[12] M\_\allowbreak{}Pl beta} & True & 20.3 \\
21 & \texttt{brane plasma at T = M8, H = M8: fraction (Pi\textasciicircum{}2 g*/\allowbreak{}30) M8\textasciicircum{}4/\allowbreak{}(3 M\_\allowbreak{}Pl\textasciicircum{}2 M8\textasciicircum{}2) < 10\textasciicircum{}-27} & True & 21.4 \\
22 & \texttt{through brane fields alone a 10 GeV radion would decay far too late: T\_\allowbreak{}RH < 10\textasciicircum{}-6 GeV (nucleosynthesis needs 4 10\textasciicircum{}-3 GeV); 10\textasciicircum{}13 Kaluza-Klein channels below m\_\allowbreak{}r/\allowbreak{}2} & True & 22.2 \\
23 & \texttt{the radion decays into the bulk: the lattice sum over the Kaluza-Klein modes of dirac16complex (m << m\_\allowbreak{}r/\allowbreak{}2) equals its continuum limit to 10\textasciicircum{}-6, the number of momenta to 10\textasciicircum{}-5, and Gamma\_\allowbreak{}bulk/\allowbreak{}Gamma\_\allowbreak{}brane > 10\textasciicircum{}12} & True & 22.2 \\
24 & \texttt{g* limits: 106.75 above the electroweak scale (to 0.1), 10.75 near 3 MeV (to 0.1), 2 + (21/\allowbreak{}4)(4/\allowbreak{}11)\textasciicircum{}(4/\allowbreak{}3)(3.044/\allowbreak{}3) = 3.3855 after e+e- annihilation (to 10\textasciicircum{}-3)} & True & 22.3 \\
25 & \texttt{radiation era: t(1 MeV) = 0.74 s within 5\% (1/\allowbreak{}(2 Ha) with g* = 10.75)} & True & 22.3 \\
26 & \texttt{age of the universe from the brane Friedmann equation: 13.79 Gyr (Planck 2018: 13.787 +- 0.020)} & True & 22.4 \\
27 & \texttt{far future: Ha(t) = H\_\allowbreak{}inf Coth[3 H\_\allowbreak{}inf t/\allowbreak{}2] agrees with the full brane solution at a = 3 and a = 10 to 10\textasciicircum{}-4} & True & 23.1 \\
28 & \texttt{couplings of the lightest mode: with V\_\allowbreak{}perp = M\_\allowbreak{}Pl\textasciicircum{}2/\allowbreak{}M8\textasciicircum{}6, lambda\_\allowbreak{}P/\allowbreak{}V\_\allowbreak{}perp == c\_\allowbreak{}P M8/\allowbreak{}M\_\allowbreak{}Pl\textasciicircum{}2, y/\allowbreak{}Sqrt[V\_\allowbreak{}perp] == c\_\allowbreak{}y M8/\allowbreak{}M\_\allowbreak{}Pl, G\_\allowbreak{}P/\allowbreak{}V\_\allowbreak{}perp == c\_\allowbreak{}V/\allowbreak{}M\_\allowbreak{}Pl\textasciicircum{}2, g\_\allowbreak{}r/\allowbreak{}V\_\allowbreak{}perp == c\_\allowbreak{}r/\allowbreak{}M\_\allowbreak{}Pl\textasciicircum{}2} & True & 24.3 \\
29 & \texttt{Boltzmann factor of production below T\_\allowbreak{}RH: Int\_\allowbreak{}0\textasciicircum{}Infinity t\textasciicircum{}3 K\_\allowbreak{}1(t) dt == 3 Pi/\allowbreak{}2 (exact), S(0) = 1, S(m\_\allowbreak{}h/\allowbreak{}T*) < 10\textasciicircum{}-2 with T* = (M8\textasciicircum{}6/\allowbreak{}M\_\allowbreak{}Pl)\textasciicircum{}(1/\allowbreak{}5) = 13.3 GeV} & True & 24.4 \\*
\end{longtable}}
Sections 7.5, 12 and 13 rest on these checks of section 26 of 7+1.nb (which contains 33 checks, 33 True; the independent program tools/ref\_radion.py has 31 boolean checks, 31 True, and the 15 groups of numbers printed by both agree to a largest relative difference of $8.2\times10^{-11}$):

{\footnotesize
\begin{longtable}{>{\raggedright\arraybackslash}p{\dimexpr 0.0500\linewidth-2\tabcolsep\relax}>{\raggedright\arraybackslash}p{\dimexpr 0.7700\linewidth-2\tabcolsep\relax}>{\raggedright\arraybackslash}p{\dimexpr 0.0800\linewidth-2\tabcolsep\relax}>{\raggedright\arraybackslash}p{\dimexpr 0.1000\linewidth-2\tabcolsep\relax}}
\toprule
\textbf{\#} & \textbf{check} & \textbf{result} & \textbf{section} \\
\midrule\endfirsthead
\toprule
\textbf{\#} & \textbf{check} & \textbf{result} & \textbf{section} \\
\midrule\endhead
\endfoot
\bottomrule\endlastfoot
1 & \texttt{the 4-sphere of radius r0: Ricci\_\allowbreak{}ij = 3 gamma\_\allowbreak{}ij (gamma the unit-sphere metric), Ricci\_\allowbreak{}ij = 0 for i != j, so R(S\textasciicircum{}4) = 12/\allowbreak{}r0\textasciicircum{}2 (in the signature of 7+1.nb the space-like block gives R = -12/\allowbreak{}r0\textasciicircum{}2, also with curvature[] of section 11); the volume of the unit S\textasciicircum{}4 is 8 Pi\textasciicircum{}2/\allowbreak{}3} & True & 26.1 \\
2 & \texttt{radion on S\textasciicircum{}4: the same reduction for ds\textasciicircum{}2 = e\textasciicircum{}(-4 beta) g\_\allowbreak{}E - e\textasciicircum{}(2 beta) r0\textasciicircum{}2 gamma (unit-sphere gamma) holds with the additional potential -(M\_\allowbreak{}Pl\textasciicircum{}2/\allowbreak{}2) e\textasciicircum{}(-4 beta) 12/\allowbreak{}(r0\textasciicircum{}2 e\textasciicircum{}(2 beta)) = -6 M\_\allowbreak{}Pl\textasciicircum{}2 e\textasciicircum{}(-6 beta)/\allowbreak{}r0\textasciicircum{}2 from the curvature of the sphere, V0 = (8 Pi\textasciicircum{}2/\allowbreak{}3) r0\textasciicircum{}4; R is independent of the angles} & True & 26.1 \\
3 & \texttt{radion on S\textasciicircum{}4 with U: Lambda + kappa U0 == 6/\allowbreak{}r0\textasciicircum{}2, kappa U1 == -12/\allowbreak{}r0\textasciicircum{}2 and m\_\allowbreak{}r\textasciicircum{}2 == kappa K\_\allowbreak{}U/\allowbreak{}12 - 2/\allowbreak{}r0\textasciicircum{}2 from V(beta); the eight-dimensional radion equation gives the same static solution and the same mass} & True & 26.2 \\
4 & \texttt{S\textasciicircum{}4 with Lambda only: dS4 x S\textasciicircum{}4 with r0\textasciicircum{}2 = 9/\allowbreak{}Lambda, H\textasciicircum{}2 = 1/\allowbreak{}r0\textasciicircum{}2 = Lambda/\allowbreak{}9 solves the 00, 11 and 44 equations; it is a MAXIMUM of V: m\textasciicircum{}2 = -6 H\textasciicircum{}2 from V(beta), and m\textasciicircum{}2 = -6/\allowbreak{}r0\textasciicircum{}2 = -6 H\textasciicircum{}2 from the eight-dimensional radion equation} & True & 26.2 \\
5 & \texttt{the radion couples to a mass M on the brane as (phi/\allowbreak{}(Sqrt[3] M\_\allowbreak{}Pl)) M: -M e\textasciicircum{}(-n beta/\allowbreak{}2) = -M + (n/\allowbreak{}2) M beta + ..., with beta = phi/\allowbreak{}(M\_\allowbreak{}Pl Sqrt[n (n + 2)/\allowbreak{}2]) the coupling is Sqrt[n/\allowbreak{}(2 (n + 2))]/\allowbreak{}M\_\allowbreak{}Pl = 1/\allowbreak{}(Sqrt[3] M\_\allowbreak{}Pl) for n = 4, and the fifth-force strength 2 (coupling M\_\allowbreak{}Pl)\textasciicircum{}2 = n/\allowbreak{}(n + 2) = 2/\allowbreak{}3} & True & 26.3 \\
6 & \texttt{zonal harmonics on S\textasciicircum{}4 (l = 0..6, exact): (1/\allowbreak{}sin\textasciicircum{}3 x4) d/\allowbreak{}dx4 (sin\textasciicircum{}3 x4 d/\allowbreak{}dx4 C\_\allowbreak{}l\textasciicircum{}(3/\allowbreak{}2)(cos x4)) == -l (l + 3) C\_\allowbreak{}l\textasciicircum{}(3/\allowbreak{}2), C\_\allowbreak{}l\textasciicircum{}(3/\allowbreak{}2)(1) == (l + 1)(l + 2)/\allowbreak{}2, and Int over S\textasciicircum{}4 of (C\_\allowbreak{}l/\allowbreak{}C\_\allowbreak{}l(1))\textasciicircum{}2 == V/\allowbreak{}d\_\allowbreak{}l (unit sphere), so the normalized zonal harmonic is Sqrt[d\_\allowbreak{}l/\allowbreak{}V] at the north pole; d\_\allowbreak{}l == Binomial[l + 4, 4] - Binomial[l + 2, 4] for l = 0..20} & True & 26.4 \\
7 & \texttt{the brane at the north pole of S\textasciicircum{}4: the embedding (r0 sin x4 sin x5 sin x6 cos x7, ..., r0 cos x4) of S\textasciicircum{}4 in R\textasciicircum{}5 induces r0\textasciicircum{}2 gamma; the north pole x4 = 0 is (0, 0, 0, 0, r0); of the 10 rotations L\_\allowbreak{}ab of R\textasciicircum{}5 (the Killing vectors of S\textasciicircum{}4, SO(5)) the 6 with a, b <= 4 vanish there (SO(4), kept by the brane) and the 4 L\_\allowbreak{}a5 move it} & True & 26.4 \\
8 & \texttt{Killing spinors of S\textasciicircum{}4 (radius r0, frame e\textasciicircum{}a = (r0, r0 sin x4, ...) dx\textasciicircum{}(a+3), gt\_\allowbreak{}a = I rho\_\allowbreak{}(a+3)): eps\_\allowbreak{}s = exp(s (i/\allowbreak{}2) x4 gt\_\allowbreak{}1) exp((1/\allowbreak{}2) x5 gt\_\allowbreak{}1 gt\_\allowbreak{}2) exp((1/\allowbreak{}2) x6 gt\_\allowbreak{}2 gt\_\allowbreak{}3) exp((1/\allowbreak{}2) x7 gt\_\allowbreak{}3 gt\_\allowbreak{}4) eps0 satisfies nabla\_\allowbreak{}mu eps\_\allowbreak{}s == s (i/\allowbreak{}(2 r0)) e\_\allowbreak{}mu\textasciicircum{}a gt\_\allowbreak{}a eps\_\allowbreak{}s (s = +-1, all eps0), so the Dirac operator gives (2 s i/\allowbreak{}r0) eps\_\allowbreak{}s: |lambda| = 2/\allowbreak{}r0, Friedrich's bound n R/\allowbreak{}(4(n - 1)) = 4/\allowbreak{}r0\textasciicircum{}2} & True & 26.4 \\
9 & \texttt{Dirac spectrum on S\textasciicircum{}4: Weyl's law for spinors (4 components) N(|lambda| <= L/\allowbreak{}r) \textasciitilde{} 4 V (L/\allowbreak{}r)\textasciicircum{}4/\allowbreak{}(32 Pi\textasciicircum{}2) = L\textasciicircum{}4/\allowbreak{}3 agrees with the degeneracies 2 x 4 Binomial[l + 3, 3] of the eigenvalues +-i(l + 2)/\allowbreak{}r (to 2.0 10\textasciicircum{}-3 at L = 1000)} & True & 26.4 \\
10 & \texttt{the sphere with the volume of the torus today (V = M\_\allowbreak{}Pl\textasciicircum{}2/\allowbreak{}M8\textasciicircum{}6): r = (3 V/\allowbreak{}(8 Pi\textasciicircum{}2))\textasciicircum{}(1/\allowbreak{}4) = 688.9 GeV\textasciicircum{}-1 = 1.359 10\textasciicircum{}-13 m; lowest scalar levels 4.03 MeV (torus, 8 states) and 2.90 MeV (sphere, l = 1, 5 states); lowest dirac16complex level on the sphere 2/\allowbreak{}r = 2.90 MeV (8 Dirac fermions; the torus has 4 massless ones for m = 0); levels below 5 GeV: torus 1.17335 10\textasciicircum{}13, sphere 1.17375 10\textasciicircum{}13, Weyl 1.17335 10\textasciicircum{}13} & True & 26.4 \\
11 & \texttt{RS1 in five dimensions (mostly minus): G\textasciicircum{}mu\_\allowbreak{}nu = 3 A'' - 6 A'\textasciicircum{}2, G\textasciicircum{}y\_\allowbreak{}y = -6 A'\textasciicircum{}2; bulk A = k y needs Lambda = -6 k\textasciicircum{}2; the delta functions A'' = +-2 k delta at the branes need the tensions f\_\allowbreak{}UV = +6 k/\allowbreak{}kappa, f\_\allowbreak{}IR = -6 k/\allowbreak{}kappa (UV positive)} & True & 26.5 \\
12 & \texttt{Israel junction at the UV brane of RS1: with K\textasciicircum{}mu\_\allowbreak{}nu = (1/\allowbreak{}2) g\textasciicircum{}mu alpha d\_\allowbreak{}y g\_\allowbreak{}alpha nu = -A' delta\textasciicircum{}mu\_\allowbreak{}nu (normal d\_\allowbreak{}y), [K\textasciicircum{}mu\_\allowbreak{}nu] - delta\textasciicircum{}mu\_\allowbreak{}nu [K] = 6 k = (coefficient of A'' in G\textasciicircum{}mu\_\allowbreak{}nu) x (jump of A'), the delta-function part of G\textasciicircum{}mu\_\allowbreak{}nu; in the signature (-,+,+,+,+) of the note G\textasciicircum{}A\_\allowbreak{}B changes sign, so there [K\_\allowbreak{}mu nu - g\_\allowbreak{}mu nu K] = -kappa S\_\allowbreak{}mu nu (the note uses +kappa S\_\allowbreak{}mu nu, hence its reversed tensions)} & True & 26.5 \\
13 & \texttt{RS1 radion: with R promoted to T(x), the bulk action is (3/\allowbreak{}(kappa k)) (d e\textasciicircum{}(-k Pi T))\textasciicircum{}2 + V(T) + (a total derivative), so the canonical radion is mu = Sqrt[6/\allowbreak{}(kappa k)] e\textasciicircum{}(-k Pi T) (the note: Sqrt[12/\allowbreak{}(kappa k)])} & True & 26.5 \\
14 & \texttt{eight dimensions, warped along one extra direction (ds\textasciicircum{}2 = e\textasciicircum{}(-2 A(y)) eta - dy\textasciicircum{}2 - dz.dz): G\textasciicircum{}mu\_\allowbreak{}nu = 3 A'' - 6 A'\textasciicircum{}2, G\textasciicircum{}y\_\allowbreak{}y = -6 A'\textasciicircum{}2, G\textasciicircum{}z\_\allowbreak{}z = 4 A'' - 10 A'\textasciicircum{}2; with only Lambda (all G\textasciicircum{}A\_\allowbreak{}A equal) A'' = 0 and A' = 0: the note's five-dimensional equations do not carry over} & True & 26.5 \\
15 & \texttt{Goldberger-Wise radion mass: V(mu) = 4 k (mu/\allowbreak{}f)\textasciicircum{}4 (v\_\allowbreak{}B - v\_\allowbreak{}A (mu/\allowbreak{}f)\textasciicircum{}eps)\textasciicircum{}2 has V'(<mu>) = 0 at (<mu>/\allowbreak{}f)\textasciicircum{}eps = v\_\allowbreak{}B/\allowbreak{}v\_\allowbreak{}A (k Pi <T> = (4 k\textasciicircum{}2/\allowbreak{}m\textasciicircum{}2) Log[v\_\allowbreak{}A/\allowbreak{}v\_\allowbreak{}B]) and V''(<mu>) = 8 k eps\textasciicircum{}2 v\_\allowbreak{}B\textasciicircum{}2 <mu>\textasciicircum{}2/\allowbreak{}f\textasciicircum{}4; with f\textasciicircum{}2 = 6/\allowbreak{}(kappa k): m\_\allowbreak{}r\textasciicircum{}2 = (4/\allowbreak{}3) kappa k\textasciicircum{}2 eps\textasciicircum{}2 v\_\allowbreak{}B\textasciicircum{}2 (<mu>/\allowbreak{}f)\textasciicircum{}2 (the note's f\textasciicircum{}2 = 12/\allowbreak{}(kappa k) gives 2/\allowbreak{}3)} & True & 26.5 \\
16 & \texttt{conformally flat brane metric h = Omega\textasciicircum{}2 eta (four dimensions, Omega(x0, x1)): Sqrt|h| R(h) == -6 Omega box Omega with box = d0\textasciicircum{}2 - d1\textasciicircum{}2 (curvature conventions of 7+1.nb)} & True & 26.6 \\
17 & \texttt{the Higgs boson on the brane (h = Omega\textasciicircum{}2 eta, Omega = 1 + r/\allowbreak{}Lambda + a2 r\textasciicircum{}2, phi\_\allowbreak{}H = v + h, gamma = v/\allowbreak{}Lambda): the terms linear in (h, r) are total derivatives, and to second order L == (1/\allowbreak{}2)(dh)\textasciicircum{}2 + (1/\allowbreak{}2)(1 + 6 xi gamma\textasciicircum{}2)(dr)\textasciicircum{}2 + 6 xi gamma dh.dr - (1/\allowbreak{}2) m\_\allowbreak{}h0\textasciicircum{}2 h\textasciicircum{}2 - (1/\allowbreak{}2) m\_\allowbreak{}r0\textasciicircum{}2 r\textasciicircum{}2 up to total derivatives (Euler-Lagrange expressions in x0 and x1); a2 drops out} & True & 26.6 \\
18 & \texttt{canonical kinetic terms: Transpose[S].K.S == 1 for S = \{\{1, -6 xi gamma/\allowbreak{}Z\}, \{0, 1/\allowbreak{}Z\}\} if and only if Z\textasciicircum{}2 = 1 + 6 xi gamma\textasciicircum{}2 (1 - 6 xi) (the note's Z); with the note's literal kinetic terms (no 6 xi gamma\textasciicircum{}2 for the radion) it would be Z\textasciicircum{}2 = 1 - 36 xi\textasciicircum{}2 gamma\textasciicircum{}2} & True & 26.6 \\
19 & \texttt{mass mixing: with (h', r') = Q(theta).(h\_\allowbreak{}m, r\_\allowbreak{}m), Q = \{\{cos, -sin\}, \{sin, cos\}\}, the off-diagonal element of Transpose[Q].M'.Q is -(1/\allowbreak{}2)(M'\_\allowbreak{}11 - M'\_\allowbreak{}22) sin 2 theta + M'\_\allowbreak{}12 cos 2 theta, zero for tan 2 theta = 2 M'\_\allowbreak{}12/\allowbreak{}(M'\_\allowbreak{}11 - M'\_\allowbreak{}22) = 12 gamma xi Z m\_\allowbreak{}h0\textasciicircum{}2/\allowbreak{}(m\_\allowbreak{}r0\textasciicircum{}2 - m\_\allowbreak{}h0\textasciicircum{}2 (Z\textasciicircum{}2 - 36 xi\textasciicircum{}2 gamma\textasciicircum{}2)) (the note's formula); with the inverse matrix (the note's display) it is +(1/\allowbreak{}2)(M'\_\allowbreak{}11 - M'\_\allowbreak{}22) sin 2 theta + M'\_\allowbreak{}12 cos 2 theta (theta -> -theta); (h, r) = S.Q.(h\_\allowbreak{}m, r\_\allowbreak{}m) gives C\_\allowbreak{}h = cos theta - (6 xi gamma/\allowbreak{}Z) sin theta, C\_\allowbreak{}r = sin theta/\allowbreak{}Z} & True & 26.6 \\
20 & \texttt{Higgs-radion mixing in 7+1 (gamma = v/\allowbreak{}(Sqrt[3] M\_\allowbreak{}Pl), xi = 1, m\_\allowbreak{}r0 = 10 GeV, m\_\allowbreak{}h0 = 125.20 GeV; 60-digit arithmetic): |Z - 1| < 10\textasciicircum{}-30, |theta| = 3.5 10\textasciicircum{}-16, the rates of h -> gamma gamma and h -> g g change by less than 10\textasciicircum{}-30; in RS1 with <mu> = 3 TeV gamma = v/\allowbreak{}(3 TeV) = 0.082} & True & 26.6 \\*
\end{longtable}}
\Needspace{12\baselineskip}\section*{16. How to reproduce}\addcontentsline{toc}{section}{16. How to reproduce}
\begin{enumerate}
\item Write the unevaluated notebook with \texttt{python3 tools/\allowbreak{}make\_\allowbreak{}7p1\_\allowbreak{}nb.py}; sections 25 and 26 are generated by \texttt{tools/\allowbreak{}nb\_\allowbreak{}part\_\allowbreak{}torus\_\allowbreak{}brane.py} and \texttt{tools/\allowbreak{}nb\_\allowbreak{}part\_\allowbreak{}radion.py}.
\item Evaluate 7+1.nb in a fresh kernel of Wolfram 15.0.1 (Evaluation \texorpdfstring{\ensuremath{\triangleright}}{▸} Quit Kernel, then Evaluation \texorpdfstring{\ensuremath{\triangleright}}{▸} Evaluate Notebook; or \texttt{Get["D:/\allowbreak{}Developer/\allowbreak{}github/\allowbreak{}7+1/\allowbreak{}tools/\allowbreak{}run\_\allowbreak{}eval.wl"]}) and save it.
\item Run \texttt{python3 tools/\allowbreak{}ref\_\allowbreak{}torus\_\allowbreak{}brane.py} and \texttt{python3 tools/\allowbreak{}ref\_\allowbreak{}radion.py}: they repeat the examples independently (exit code 0 if all pass).
\item Run \texttt{python3 tools/\allowbreak{}make\_\allowbreak{}docs.py}: it regenerates this document, its symbol dictionary and its provenance record (and all other documents).
\end{enumerate}
\Needspace{12\baselineskip}\section*{17. References}\addcontentsline{toc}{section}{17. References}
\begin{itemize}
\item K. Akama, Lect. Notes Phys. 176 (1982) 267; V. A. Rubakov, M. E. Shaposhnikov, Phys. Lett. B 125 (1983) 136 — matter confined to a submanifold of a higher-dimensional spacetime.
\item N. Arkani-Hamed, S. Dimopoulos, G. Dvali, Phys. Lett. B 429 (1998) 263 and Phys. Rev. D 59 (1999) 086004 — the Standard Model on a brane with large extra dimensions; the normalcy temperature.
\item L. Randall, R. Sundrum, Phys. Rev. Lett. 83 (1999) 3370 — branes in a warped bulk (section 13; the bulk of 7+1 is not warped); W. D. Goldberger, M. B. Wise, Phys. Rev. Lett. 83 (1999) 4922 and Phys. Lett. B 475 (2000) 275 — the stabilized warped radion; W. Israel, Nuovo Cim. B 44 (1966) 1 — junction conditions.
\item G. F. Giudice, R. Rattazzi, J. D. Wells, Nucl. Phys. B 595 (2001) 250 — the curvature coupling of the Higgs boson and Higgs–radion mixing.
\item R. Camporesi, A. Higuchi, J. Geom. Phys. 20 (1996) 1 — the Dirac spectrum on spheres; Th. Friedrich, Math. Nachr. 97 (1980) 117 — its lower bound.
\item G. Dvali, G. Gabadadze, M. Porrati, Phys. Lett. B 485 (2000) 208 — brane-localized curvature (not part of this model).
\item R. Sundrum, Phys. Rev. D 59 (1999) 085009 — brane fluctuations (branons; not part of this model).
\item S. Navas et al. (Particle Data Group), Phys. Rev. D 110 (2024) 030001; Planck Collaboration, Astron. Astrophys. 641 (2020) A6.
\end{itemize}
\end{document}
